% ============================================================
%  第九卷 第二章（后半） T=0 时费米系统的格林函数（§14–§21）
%  底本：高教社中文第4版（王锡绂 译）；印刷页 49–77（PDF p060–p088）
%  不带 \chapter（章头在 ch02a.tex），自 \section{自能函数} 起
%
%  插图清单：本章无编号插图（图 1 在前半章）。原书"图式等式"（带公式编号
%  的费曼图公式）共 15 处：(14.3)(14.4)(14.8)(14.9)(14.10)(15.9)(15.10)
%  (15.11)(16.1)(17.2)(17.3)(17.4)(19.3)(21.1)(21.6)，另有无编号图式 7 处
%  （p064 零级双图、p065 两处、p068 交叉虚线图、p071 级数图、p074 危险图、
%  p077 外虚线元素），全部用 picture 环境按原书页面重绘（外线标注、箭头
%  方向逐图对照 p060–p088 页面 PNG）。
%
%  圈码脚注：22 处已全部转录为 \footnote{...}（p061、p062、p063、p065、p066、
%  p067、p068、p070×2、p074×2、p075、p076、p077、p078、p080×2、p083、p084、
%  p086×2、p088），脚注文本按 原书转换/pages 相应页面 300dpi 渲染逐字转录，
%  原书脚注按页重编号（每页自①起），与 \footnote 按页自动复位一致。
%
%  编号公式核算：OCR 带 tag 的 98 条（其中 2 条乱码已按页面订正：
%  "5.14"→(15.14)（p067）、".17"→(21.17)（p087)）。对照页面核实，OCR 有
%  11 处图式等式编号未带 tag（编号见于 OCR 正文文本）：(14.4)(14.8)(14.9)
%  (14.10)(15.9)(15.11)(16.1)(17.2)(17.3)(17.4)(21.1)；另有纯公式 2 处漏 tag：
%  (21.13)（p086）、(21.16)（p087），页面均带编号。故本章编号公式共
%  98+13=111 个；逐节与页面核对无重号、无缺号、无跳号。
%
%  已修正的 OCR 错误（均对照 原书转换/pages/p0NN.png 逐条核实）：
%   1. §14 首段（p060/p061）："与其出现在什样的图中无关"→"什么样"（p061）；
%      "我们以 iΣ_αβ(P) 表示它"OCR 丢负号，按 p061 为 −iΣ_αβ(P)。
%   2. (14.3) 后正文：OCR "$(\mathbf{\Gamma}-\mathrm{i}\Sigma_{\alpha\beta})$"
%      按 p061 为 (−iΣ_αβ)。
%   3. §14（p062）：OCR "分枝"图之和记号 "$\mathscr{L}_a$" 按 p062 为 Σ_a；
%      (14.12) OCR n(**μ**) 按 p062 为 n(μ)（μ 不加粗）。
%   4. §15（p063）："需要重新改用相互作用绘景中的ψ算符.ψ如讨论函数_G时"
%      按 p063 接排为"…ψ算符.如讨论函数 G 时的情况一样…"。
%   5. (15.5)（p064）：OCR 指数重复乱串，按 p064 重建为
%      exp{−i[P₃(X₃−X₂)+P₄(X₄−X₂)−P₁(X₁−X₂)]}。
%   6. (15.6)–(15.8)（p064）：OCR "\hat{\theta}"/"\mathbb{\Lambda}^{(4)}"
%      按 p064 均为 δ^{(4)}；花体 \cal G 按原书为斜体 G。
%   7. (15.12)（p066）：OCR 丢失首行左端，按 p066 补
%      "(i∂/∂t₁+Δ₁/2m+μ)G_αβ(X₁−X₂)−δ_αβδ^{(4)}(X₁−X₂)="。
%   8. §16（p068）：OCR "从数学观ω点来说"→"从数学观点来说"；
%      "散射辐"→"散射幅"（p068/p069/p086 多处，页面均作"幅"）；
%      (16.2) 上标 OCR 乱码按 p068 为 G^{(真空)}。
%   9. §17（p070/p071）：OCR 倒装"因此Γ只Γ̃是变量"按 p070 接排；
%      (17.5) 下标按 p071 为 Γ̃_{γζ,ακ}、Γ_{κδ,ζβ}（OCR γδ,ακ/γβ 之误）；
%      (17.7) OCR "]^{\top}"按 p071 为 ]^{−1}；(17.10) OCR "\hat 8"为 δ、
%      v_r 为 v_F、𝒍·𝒌 加粗。
%  10. (17.12)(17.16)（p073）：极限后置语 OCR 乱码按 p073 为"，当 k/ω→0"、
%      "，当 ω/k→0"；(17.15) OCR "γξ,ακ"、"dω_l"按 p073 为 γζ,ακ、do_l。
%  11. §17 末（p074）：OCR "P↑−P↑≡S"按 p074 为 P₁−P₂≡S=(s₀,s)。
%  12. §18（p075）：OCR "当q+k p_F p_ξ基态激发时"按 p075 接排为
%      "而另一个应小于 p_F.从另一方面来看，当基态激发时"；(18.1) 宗量
%      (n,t)→(n,l)、χ_{κζ}(t)→χ_{κζ}(l)、"doσ_ι"→do_l（p076）；
%      (18.3) v_p→v_F、"∫dzΓ"→∫Γ、do′（p076）。
%  13. §19（p077/p078）：(19.1) OCR "\hat\otimes"按 p077 为 δ（δV̂^{(1)}、
%      δÛ）；(19.2) 指数按 p077 为 e^{iP₂X}…e^{−iP₁X}；p078 "算符 ψ*(X)
%      和 Ψ(X)"按 p078 为 Ψ̂⁺(X) 和 Ψ̂(X)；(19.6) "\hat\otimes^{(4)}"为
%      δ^{(4)}；p078 末 Ψ̂′⁺_α 的星号按 p078 为加号。
%  14. (19.10)（p079）：OCR "YYω→0 时"按 p079 为"，当 k/ω→0 时"。
%  15. §19（p080）：OCR "μ+8U=常量"→μ+δU=常量；δG 傅里叶分量式
%      "\mathfrak F^{(4)}"→δ^{(4)}、"\dot U_0"→U₀；(19.15) 下标
%      "\delta_{\infty\beta}"→δ_{αβ}。
%  16. §20（p081/p082）：OCR 脱文"以下计算的目的，是将等式的右边部分只用沿"
%      与"首先，用(17.17)的表达式(其中将记号 Q_F 换成 S_F)代替上面的第二个
%      积分中的 Γ^k；"按 p081 补齐；(20.3) 下标按 p082 为 Γ^ω_{αζ,ακ}、
%      Γ^k_{κγ,ζγ}；(19.17) OCR "dt dt dt"乱码按 p081 重建。
%  17. §20（p082/p083）：OCR 残句"然后对dos求积分又得出因子4π.因而…"按
%      p082 接排；p083 δn′ 式按页面重建为 δn′=(∂n′/∂p_F)δp_F=δ(p−p_F)δp_F；
%      "[指(2.14)一译注]"→"—译注"。
%  18. §21（p083/p084）：OCR "[B. M. Γ，1958]"按 p083 为 [В. М. Галицкий,
%      1958]（西里尔）；p084 "$\hat\psi$ 和 $\hat\psi^\star$"按页面为
%      Ψ̂ 和 Ψ̂⁺（大写）；"[Σ]_a"的 OCR 乱码下标按 p084 为 a；(21.3) 后
%      "显然比 PF 大得多"→p_F。
%  19. (21.11)（p085）：OCR "p_ℓ+p_ℓ"按 p085 为 |p₁+p₂−p′|。
%  20. (21.12) 后代换式（p085）：OCR "{\right.}…{\left.}"按 p085 重建为
%      sign(p′−p_F)→1−θ(p′)−θ(p₁+p₂−p′)。
%  21. §21（p086）：OCR "G_{\mathfrak o}(Q)"按 p086 为 G₀(Q)；
%      "−\operatorname{Re}f^{\mathcal Q}"按 p086 为"−Re f②"。
%  22. (21.17)（p087）：OCR 尾部乱码"(21−dθ(p′))"系噪声，按 p087 该式至
%      d³q d³p′/(2π)⁶ 止（无标点）。(21.18) 后"(此处…只要取 ε(p)=p²/2m"
%      按 p087 补"就够了).按照上面指出的意义，以"。
%  23. 全章：OCR 全角句号"．"统一为半角"."；\pmb 一律改 \bm；
%      函数名收紧（\exp,\operatorname{sign},\operatorname{Re} 等）；
%      p070 末"AH"、p070"22"、p071 后"ANU"等 OCR 噪声删除。
%
%  原书排印问题（照排并登记，供总校裁定）：
%   a. (18.7) 第二式（p076）：A 的展开按页面照排为 "B(ϑ)δ_αγδ_βδ+
%      G(ϑ)σ_γα·σ_δβ"——第二个系数页面印作 G，按上下文（"其中系数
%      F,G,B,C"及 (18.9) 的 C_l=G_l(1−C_l)）应为 C，疑排印之误。
%   b. (21.13)（p086）：分子按页面照排为 θ(p₁+p₂+p′)，按 (21.12) 的代换
%      及 (21.16) 的对应项 θ(p+q−p′) 应为减号，疑排印之误（+→−）。
%   c. p080（伽利略变换段）："出现在函数 G(F) 中"按页面照排，按上下文
%      应为 G(P)。
%   d. p080 (19.15)：按页面右端无闭花括号（至 d⁴Q/(2π)⁴ 后直接句号），
%      照排不补。
%   e. p080 (19.14) 后 δU(P₂,P₁) 式中动量下标按页面照排为 p_i
%      （−p_i·w₀…），疑为 p₁ 之误。
%   f. p075 "其中 n 为变单位矢量"按页面照排（"变"字疑衍）。
%
%  遗留问题：
%   - 图式等式为按页面重绘的示意图（无 TikZ，用 picture 环境），拓扑与
%     外线标注与原书一致，内线箭头等细节或有简省；
%   - (14.9)(14.10)(15.11)(21.1) 多图组的 (a)(b)(c)… 编号按原书置于图下。
% ============================================================

% ---- 图式等式辅助记号（picture 环境内使用） ----
\newcommand{\arrL}[3]{\put(#1,#2){\vector(-1,0){#3}}}   % 细实线，左向箭头
\newcommand{\tarrL}[3]{\put(#1,#2){{\thicklines\vector(-1,0){#3}}}} % 粗实线，左向箭头
\newcommand{\dvh}[3]{\multiput(#1,#2)(0,7){#3}{\line(0,1){5}}} % 竖虚线（节距7节长5）

\section{自能函数}

上一节表述的图技术规则具有重要性质：图中的公共系数与图的级无关.由于这个性质，图中每个“图元素”都具有确定的解析意义，而与其出现在什么样的图中无关，因此图元素可以预先独立地计算出来.并且可以预先计算出某些具有一定数量端点的图元素之和，然后再把这个“单元”组装入更复杂的图中.这是图技术最重要的优点之一.

有一种“单元”，它也有重要的独立意义，即所谓的自能函数\footnote{对照量子电动力学中类似的定义，那里将这个函数称作紧致自能函数(第四卷 §103，§105).}.为了得出这个概念，现在来分析一下格林函数图中一切不能靠切断一条实联线而分为两部分的图形.例如，一级微扰论的两个图形(13.13)和二级微扰论的图形(13.14a—e)都属于这种图形.所有这些图形都有同样的构造：每个端点有一个因子 $\mathrm{i}G^{(0)}_{\alpha\beta}$，还有一个称作自能函数的内部部件($P$ 的函数).所有可能的内部部件之和，称作精确的或完全的自能函数，又称质量算符；我们以 $-\mathrm{i}\Sigma_{\alpha\beta}(P)$ 表示它.

自能函数的全部图形对格林函数的贡献等于
\begin{equation}
  \setlength{\unitlength}{0.85pt}%
  \begin{picture}(480,26)(0,0)
    \tarrL{58}{12}{52}%
    \put(74,12){\makebox(0,0){$=$}}%
    \arrL{126}{12}{40}%
    \put(140,12){\makebox(0,0){$+$}}%
    \arrL{196}{12}{36}%
    \put(202,12){\circle{14}}%
    \arrL{246}{12}{36}%
    \put(258,12){\makebox(0,0){$+$}}%
    \arrL{310}{12}{36}%
    \put(316,12){\circle{14}}%
    \arrL{360}{12}{36}%
    \put(366,12){\circle{14}}%
    \arrL{410}{12}{36}%
    \put(424,12){\makebox(0,0)[l]{$+\cdots,$}}%
  \end{picture}
\end{equation}
式中除 $G^{(0)}_{\alpha\beta}=G^{(0)}\delta_{\alpha\beta}$ 之外，同样也可写出
\begin{equation}
  \Sigma_{\alpha\beta}(P)=\delta_{\alpha\beta}\Sigma(P).
\end{equation}
完全格林函数(图中以粗实线表示)由下列无穷级数之和给出：
\begin{equation}
  \setlength{\unitlength}{0.85pt}%
  \begin{picture}(480,26)(0,0)
    \tarrL{58}{12}{52}%
    \put(74,12){\makebox(0,0){$=$}}%
    \arrL{126}{12}{40}%
    \put(140,12){\makebox(0,0){$+$}}%
    \arrL{196}{12}{36}%
    \put(202,12){\circle{14}}%
    \arrL{246}{12}{36}%
    \put(258,12){\makebox(0,0){$+$}}%
    \arrL{310}{12}{36}%
    \put(316,12){\circle{14}}%
    \arrL{360}{12}{36}%
    \put(366,12){\circle{14}}%
    \arrL{410}{12}{36}%
    \put(424,12){\makebox(0,0)[l]{$+\cdots,$}}%
  \end{picture}
\end{equation}
图中小圆表示精确的自能函数 $(-\mathrm{i}\Sigma_{\alpha\beta})$.这个级数(从第三项起)的每一项乃是图形的集合，这些图形还能截成彼此间曾以一条实线相连的两个、三个等等部件.

如果从级数(14.3)的第二项开始将各项中的一个小圆连同其右边的一条连线“切除”，则剩下的级数重新与整个级数相等.这就是说：
\begin{equation}
  \setlength{\unitlength}{0.85pt}%
  \begin{picture}(260,26)(0,0)
    \tarrL{56}{12}{50}%
    \put(74,12){\makebox(0,0){$=$}}%
    \arrL{122}{12}{38}%
    \put(136,12){\makebox(0,0){$+$}}%
    \qbezier(178,4)(185.312,4.688)(186,12)%
    \qbezier(186,12)(185.312,19.312)(178,20)%
    \qbezier(178,20)(170.688,19.312)(170,12)%
    \qbezier(170,12)(170.688,4.688)(178,4)%
    \tarrL{222}{12}{52}%
    \arrL{242}{12}{54}%
  \end{picture}
\end{equation}
将这个等式写成解析形式，如：
\begin{equation}
  G=G^{(0)}+G\Sigma G^{(0)},
\end{equation}
或除以 $G^{(0)}G$：
\begin{equation}
  \frac{1}{G(P)}=\frac{1}{G^{(0)}(P)}-\Sigma(P).
\end{equation}
我们指出，$\Sigma$ 的虚部的符号与 $\operatorname{Im}G$ 的符号相同，根据(8.14)式有
\begin{equation}
  \operatorname{sign}\operatorname{Im}\Sigma(\omega,\bm{p})=-\operatorname{sign}\omega.
\end{equation}
这个等式是考虑到 $\operatorname{Im}G^{-1}$ 与 $\operatorname{Im}G$ 的符号相反由(14.6)式得出的，依照(9.7)式，$\operatorname{Im}G^{(0)-1}=0$.

因此，计算 $G$ 归结为计算 $\Sigma$，这只需研究少数图形.这个图形数还可以再减少，因为剩余图形的部件会立刻相加成很简单的表达式.

就是说我们从决定 $\Sigma$(粒子间是成对相互作用)的全部图形集合中分离出各种“分枝”图，这些“分枝”图是用一条虚线连接到各外线上的：它们之和以 $\Sigma_a$ 标记.所有这些图形都包含在如下形状的一个骨架图形之中\footnote{如同量子场论，将粗线和单元组成的图形称作骨架图；每个这样的图都等价于无穷多个不同级的一般图形的一定集合.}：
\begin{equation}
  \setlength{\unitlength}{0.85pt}%
  \begin{picture}(360,130)(0,0)
    \arrL{50}{74}{46}%
    \qbezier(86,50)(107.936,52.064)(110,74)%
    \qbezier(110,74)(107.936,95.936)(86,98)%
    \qbezier(86,98)(64.064,95.936)(62,74)%
    \qbezier(62,74)(64.064,52.064)(86,50)%
    \put(86,74){\makebox(0,0){$-\mathrm{i}\Sigma_a$}}%
    \arrL{158}{74}{44}%
    \put(180,74){\makebox(0,0){$=$}}%
    \qbezier(250,76)(276,100)(250,124)%
    \qbezier(250,76)(224,100)(250,124)%
    \put(263,92){\vector(0,-1){9}}%
    \dvh{250}{40}{5}%
    \arrL{302}{38}{50}%
    \arrL{346}{38}{40}%
  \end{picture}
\end{equation}
$\Sigma$ 的其余部分用 $\Sigma_b$ 标记.这样，在一级和二级图形中，属于第一种的图形如下：
\begin{equation}
  \setlength{\unitlength}{0.66pt}%
  \begin{picture}(660,155)(0,0)
    \arrL{50}{84}{44}%
    \qbezier(76,62)(96.108,63.892)(98,84)%
    \qbezier(98,84)(96.108,104.108)(76,106)%
    \qbezier(76,106)(55.892,104.108)(54,84)%
    \qbezier(54,84)(55.892,63.892)(76,62)%
    \put(76,84){\makebox(0,0){$-\mathrm{i}\Sigma_a$}}%
    \arrL{146}{84}{42}%
    \put(166,84){\makebox(0,0){$=$}}%
    \qbezier(230,89)(252,106)(230,123)%
    \qbezier(230,89)(208,106)(230,123)%
    \put(236,95){\vector(0,-1){8}}%
    \dvh{230}{50}{5}%
    \arrL{272}{48}{40}%
    \arrL{312}{48}{38}%
    \put(230,24){\makebox(0,0){(a)}}%
    \put(330,84){\makebox(0,0){$+$}}%
    \qbezier(362,108)(395,132)(428,108)%
    \multiput(364,108)(7,0){9}{\line(1,0){5}}%
    \put(366,106){\line(1,-1){28}}%
    \put(424,106){\line(-1,-1){28}}%
    \put(404,86){\vector(-1,-1){9}}%
    \dvh{395}{50}{3}%
    \arrL{437}{48}{40}%
    \arrL{477}{48}{38}%
    \put(395,24){\makebox(0,0){(b)}}%
    \put(494,84){\makebox(0,0){$+$}}%
    \qbezier(560,118)(580,132)(560,146)%
    \qbezier(560,118)(540,132)(560,146)%
    \put(565,122){\vector(0,-1){8}}%
    \dvh{560}{102}{3}%
    \qbezier(560,100)(524,82)(560,64)%
    \qbezier(560,100)(596,82)(560,64)%
    \put(546,86){\vector(-1,-1){8}}%
    \dvh{560}{50}{2}%
    \arrL{602}{48}{38}%
    \arrL{640}{48}{36}%
    \put(560,24){\makebox(0,0){(c)}}%
  \end{picture}
\end{equation}
属于第二种的图形为：
\begin{equation}
  \setlength{\unitlength}{0.62pt}%
  \begin{picture}(700,205)(0,0)
    \arrL{54}{150}{46}%
    \qbezier(82,128)(102.108,129.892)(104,150)%
    \qbezier(104,150)(102.108,170.108)(82,172)%
    \qbezier(82,172)(61.892,170.108)(60,150)%
    \qbezier(60,150)(61.892,129.892)(82,128)%
    \put(82,150){\makebox(0,0){$-\mathrm{i}\Sigma_b$}}%
    \arrL{152}{150}{42}%
    \put(172,150){\makebox(0,0){$=$}}%
    \put(190,150){\line(1,1){9}}%
    \put(202,161){\line(1,0){8}}%
    \put(214,160){\line(1,-1){9}}%
    \arrL{278}{150}{50}%
    \arrL{324}{150}{42}%
    \put(230,118){\makebox(0,0){(a)}}%
    \put(340,150){\makebox(0,0){$+$}}%
    \put(360,150){\line(1,1){13}}%
    \put(375,164){\line(2,1){18}}%
    \put(395,174){\line(2,-1){18}}%
    \put(415,164){\line(1,-1){13}}%
    \put(388,150){\line(1,1){8}}%
    \put(398,159){\line(1,0){8}}%
    \put(408,158){\line(1,-1){8}}%
    \arrL{478}{150}{48}%
    \arrL{524}{150}{42}%
    \put(400,118){\makebox(0,0){(b)}}%
    \put(540,150){\makebox(0,0){$+$}}%
    \put(560,150){\line(1,-1){9}}%
    \put(571,140){\line(1,0){8}}%
    \put(581,141){\line(1,1){9}}%
    \qbezier(606,154)(622,166)(606,178)%
    \qbezier(606,154)(590,166)(606,178)%
    \put(611,158){\vector(0,-1){8}}%
    \arrL{650}{150}{48}%
    \arrL{696}{150}{42}%
    \put(600,118){\makebox(0,0){(c)}}%
    \put(28,80){\makebox(0,0){$+$}}%
    \put(52,80){\line(1,-1){9}}%
    \put(63,70){\line(1,0){8}}%
    \put(73,71){\line(1,1){9}}%
    \put(100,80){\line(1,1){9}}%
    \put(111,90){\line(1,0){8}}%
    \put(121,89){\line(1,-1){9}}%
    \arrL{176}{80}{50}%
    \arrL{226}{80}{46}%
    \put(130,48){\makebox(0,0){(d)}}%
    \put(248,80){\makebox(0,0){$+$}}%
    \put(312,80){\line(-1,0){66}}%
    \arrL{366}{80}{54}%
    \arrL{420}{80}{48}%
    \dvh{270}{84}{4}%
    \dvh{326}{84}{4}%
    \put(272,116){\line(1,0){20}}%
    \put(304,116){\line(1,0){20}}%
    \put(296,119){\vector(1,0){14}}%
    \put(300,48){\makebox(0,0){(e)}}%
  \end{picture}
\end{equation}
(14.8)图形中的粗圈相应于系统的精确密度 $n(\mu)$［与此类似，(13.13a)图形中的细圈相应于理想气体的密度 $n^{(0)}(\mu)$］.因此由定义(14.8)可得
\begin{equation}
  -\mathrm{i}\Sigma_a=-\mathrm{i}n(\mu)U(0),
\end{equation}
所以
\begin{equation}
  \Sigma=n(\mu)U(0)+\Sigma_b.
\end{equation}
于是需要特别计算的只有 $\Sigma_b$ 中的图形.

准粒子的色散律由方程(8.16)定义.在方程中将 $G$ 用 $\Sigma$ 表达出来，按照(14.6)式并取(9.7)式的 $G^{(0)}$，我们得到该方程的下列形式：
\begin{equation}
  \frac{1}{G^{(0)}(\varepsilon-\mu,\bm{p})}=\varepsilon(\bm{p})-\frac{p^2}{2m}=\Sigma(\varepsilon-\mu,\bm{p}).
\end{equation}
在费米球界面上，即当 $p=p_{\mathrm{F}}$ 时，准粒子能量与 $\mu$ 相等.由此可见
\begin{equation}
  \mu-\Sigma(0,p_{\mathrm{F}})=\frac{p_{\mathrm{F}}^2}{2m}.
\end{equation}
结果色散律方程(当 $p$ 的值接近 $p_{\mathrm{F}}$ 时)取如下形式：
\begin{equation}
  \varepsilon(\bm{p})-\mu=\frac{p_{\mathrm{F}}}{m}(p-p_{\mathrm{F}})
  +\Sigma(\varepsilon-\mu,p_{\mathrm{F}})-\Sigma(0,p_{\mathrm{F}}).
\end{equation}
应当强调，这里 $p_{\mathrm{F}}$ 是相互作用粒子系统边界动量的精确值.该值受关系式 $p_{\mathrm{F}}^3/3\pi^2=n$ 制约，这里 $n$ 是精确密度 $n(\mu)$，而不是(13.5)式中的近似密度 $n^{(0)}$.

\section{双粒子格林函数}

研究了四个海森伯 $\psi$ 算符的 $T$ 乘积对基态的平均\footnote{我们还利用简化记号，这里下标 $1,2,\cdots$ 标记 4 维坐标及自旋下标的集合：$X_1\alpha$, $X_2\beta$, $\cdots$(见 44 页的注解). 完整地写出来是 $K_{34,12}\equiv K_{\gamma\delta,\alpha\beta}(X_3,X_4;X_1,X_2)$}：
\begin{equation}
  K_{34,12}=\langle T\hat{\Psi}_3\hat{\Psi}_4\hat{\Psi}^{+}_1\hat{\Psi}^{+}_2\rangle,
\end{equation}
我们便得到图技术其它一些重要概念.上述函数称作双粒子格林函数［因而区别于所谓单粒子格林函数(7.9)］.

为了运用微扰论和建立图技术，需要重新改用相互作用绘景中的 $\psi$ 算符.如讨论函数 $G$ 时的情况一样，这将使得在 $T$ 乘积记号下出现因子 $\hat{S}$：
\begin{equation}
  K_{34,12}=\frac{1}{\langle\hat{S}\rangle}
  \langle T\hat{\Psi}_{03}\hat{\Psi}_{04}\hat{\Psi}^{+}_{01}\hat{\Psi}^{+}_{02}\hat{S}\rangle.
\end{equation}
在零级近似下(即当 $\hat{S}=1$ 时)该表达式可分解成用 $G^{(0)}$ 函数表示的两个收缩乘积之和：
\begin{equation}
  K^{(0)}_{34,12}=G^{(0)}_{31}G^{(0)}_{42}-G^{(0)}_{32}G^{(0)}_{41}.
\end{equation}
以下将要在动量表象中讨论以这种形式定义的双粒子格林函数的性质.

对于均匀系统，函数 $K_{34,12}$ 实际上只与四个独立自变量之差(例如 $X_3-X_2$，$X_4-X_2$，$X_1-X_2$)有关.在动量表象中，这个性质表现在：按所有变量 $X_1,\cdots,X_4$ 展开的每个傅里叶分量均含一个 $\delta$ 函数：
\begin{equation}
\begin{aligned}
  &\int K_{34,12}\exp\{\mathrm{i}(P_3X_3+P_4X_4-P_1X_1-P_2X_2)\}\mathrm{d}^4X_1\cdots\mathrm{d}^4X_4=\\
  =&(2\pi)^4\delta^{(4)}(P_3+P_4-P_1-P_2)K_{\gamma\delta,\alpha\beta}(P_3,P_4;P_1,P_2).
\end{aligned}
\end{equation}
此式不难证明，只要注意到：
\begin{equation*}
\begin{aligned}
  &P_3X_3+P_4X_4-P_1X_1-P_2X_2=\\
  =&P_3(X_3-X_2)+P_4(X_4-X_2)-P_1(X_1-X_2)-X_2(P_1+P_2-P_3-P_4),
\end{aligned}
\end{equation*}
并对 $X_3-X_2$，$X_4-X_2$，$X_1-X_2$，$X_2$ 求积分即可.我们顺便指出，傅里叶逆变换公式可写成：
\begin{equation}
\begin{aligned}
  K_{34,12}=\int &K_{\gamma\delta,\alpha\beta}(P_3,P_4;P_1,P_3+P_4-P_1)\times\\
  &\times\exp\{-\mathrm{i}[P_3(X_3-X_2)+P_4(X_4-X_2)-P_1(X_1-X_2)]\}
  \frac{\mathrm{d}^4P_1\mathrm{d}^4P_3\mathrm{d}^4P_4}{(2\pi)^{12}}.
\end{aligned}
\end{equation}
用这种形式定义的函数 $K_{\gamma\delta,\alpha\beta}(P_3,P_4;P_1,P_2)$，我们将称作动量表象中的双粒子格林函数；它的宗量以如下等式相约束：
\begin{equation*}
  P_1+P_2=P_3+P_4
\end{equation*}
在零级近似，对于上述函数［与(15.3)式对应］，有
\begin{equation}
\begin{aligned}
  &K^{(0)}_{\gamma\delta,\alpha\beta}(P_3,P_4;P_1,P_2)=\\
  =&(2\pi)^4[\delta^{(4)}(P_1-P_3)G^{(0)}_{\gamma\alpha}(P_1)G^{(0)}_{\delta\beta}(P_2)
  -\delta^{(4)}(P_1-P_4)G^{(0)}_{\gamma\beta}(P_2)G^{(0)}_{\delta\alpha}(P_1)].
\end{aligned}
\end{equation}
也就是 $K$ 归结为两个单粒子格林函数乘积之和.

在微扰论的高级近似中出现的一些项，是对这些单粒子函数的修正.但是，此外也出现一些不属于 $G$ 函数乘积的项.正是双粒子格林函数的这个部分具有独立意义.为了把它分离出来，我们将 $K$ 表示成如下形式：
\begin{equation}
\begin{aligned}
  K_{\alpha_3\alpha_4,\alpha_1\alpha_2}(P_3,&P_4;P_1,P_2)=(2\pi)^4[\delta^{(4)}(P_1-P_3)G_{\alpha_3\alpha_1}(P_1)G_{\alpha_4\alpha_2}(P_2)-\\
  -&\delta^{(4)}(P_1-P_4)G_{\alpha_3\alpha_2}(P_2)G_{\alpha_4\alpha_1}(P_1)]+\\
  +&G_{\alpha_3\beta_3}(P_3)G_{\alpha_4\beta_4}(P_4)\mathrm{i}\Gamma_{\beta_3\beta_4,\beta_1\beta_2}(P_3,P_4;P_1,P_2)G_{\beta_1\alpha_1}(P_1)G_{\beta_2\alpha_2}(P_2).
\end{aligned}
\end{equation}
以这种形式定义的函数 $\Gamma$ 称作顶角函数.

根据定义(15.1)，在空间-时间表象中的双粒子格林函数，对于交换第一对宗量或第二对宗量即交换 1 和 2 或交换 3 和 4(连同自旋下标)是反对称的.由此可得动量表象中格林函数和顶角函数类似的对称性质：
\begin{equation}
\begin{aligned}
  \Gamma_{\gamma\delta,\alpha\beta}(P_3,P_4;P_1,P_2)&=-\Gamma_{\delta\gamma,\alpha\beta}(P_4,P_3;P_1,P_2)=\\
  &=-\Gamma_{\gamma\delta,\beta\alpha}(P_3,P_4;P_2,P_1)
\end{aligned}
\end{equation}
如果研究一下展开双粒子格林函数表达式(15.2)时产生的各图之特点，则定义函数 $\Gamma$［(15.7)式的末项］时分离出的四个 $G$ 因子的意义就变得明显了.以下的讨论还假定粒子之间是成对相互作用的.

在零级近似，函数 $K$ 用下图表示：
\begin{equation*}
  \setlength{\unitlength}{0.85pt}%
  \begin{picture}(320,100)(0,0)
    \put(120,88){\makebox(0,0){$P_3=P_1$}}%
    \tarrL{160}{64}{80}%
    \put(120,40){\makebox(0,0){$P_4=P_2$}}%
    \put(260,88){\makebox(0,0){$P_4=P_1$}}%
    \tarrL{300}{64}{80}%
    \put(260,40){\makebox(0,0){$P_3=P_2$}}%
  \end{picture}
\end{equation*}
这些图对应于(15.6)式的两项.在一级微扰论中有如下类型的图\footnote{如单粒子格林函数的情况一样，定义(15.2)中的因子 $\langle\hat{S}\rangle^{-1}$ 可消去不相连的实线圈图.}：
\begin{equation*}
  \setlength{\unitlength}{0.85pt}%
  \begin{picture}(360,120)(0,0)
    \arrL{150}{96}{70}%
    \arrL{150}{74}{70}%
    \put(118,74){\line(0,-1){8}}%
    \qbezier(118,42)(128.968,43.032)(130,54)%
    \qbezier(130,54)(128.968,64.968)(118,66)%
    \qbezier(118,66)(107.032,64.968)(106,54)%
    \qbezier(106,54)(107.032,43.032)(118,42)%
    \put(126,48){\vector(0,-1){8}}%
    \put(212,84){\line(1,1){10}}%
    \put(226,96){\line(1,0){10}}%
    \put(240,94){\line(1,-1){10}}%
    \put(212,60){\line(1,-1){10}}%
    \put(226,48){\line(1,0){10}}%
    \put(240,50){\line(1,1){10}}%
    \arrL{320}{84}{80}%
    \arrL{320}{60}{80}%
  \end{picture}
\end{equation*}
它们是对(15.6)式的每个独立因子的修正.但是，此外还出现未分成两个独立部分的图：
\begin{equation}
  \setlength{\unitlength}{0.68pt}%
  \begin{picture}(600,175)(0,0)
    \put(64,120){\vector(-1,1){36}}%
    \put(64,44){\vector(-1,-1){36}}%
    \qbezier(100,56)(123.764,58.236)(126,82)%
    \qbezier(126,82)(123.764,105.764)(100,108)%
    \qbezier(100,108)(76.236,105.764)(74,82)%
    \qbezier(74,82)(76.236,58.236)(100,56)%
    \put(100,82){\makebox(0,0){$\mathrm{i}\Gamma^{(1)}$}}%
    \put(150,114){\vector(-1,-1){33}}%
    \put(150,50){\vector(-1,1){33}}%
    \put(28,160){\makebox(0,0){$P_3$}}%
    \put(162,160){\makebox(0,0){$P_1$}}%
    \put(28,6){\makebox(0,0){$P_4$}}%
    \put(162,6){\makebox(0,0){$P_2$}}%
    \put(196,82){\makebox(0,0){$=$}}%
    \arrL{310}{108}{60}%
    \arrL{360}{108}{44}%
    \dvh{310}{62}{6}%
    \arrL{310}{56}{60}%
    \arrL{360}{56}{44}%
    \put(246,118){\makebox(0,0){$P_3$}}%
    \put(372,118){\makebox(0,0){$P_1$}}%
    \put(246,42){\makebox(0,0){$P_4$}}%
    \put(372,42){\makebox(0,0){$P_2$}}%
    \put(400,82){\makebox(0,0){$+$}}%
    \arrL{516}{108}{60}%
    \arrL{566}{108}{44}%
    \dvh{516}{62}{6}%
    \arrL{516}{56}{60}%
    \arrL{566}{56}{44}%
    \put(452,118){\makebox(0,0){$P_4$}}%
    \put(578,118){\makebox(0,0){$P_1$}}%
    \put(452,42){\makebox(0,0){$P_3$}}%
    \put(578,42){\makebox(0,0){$P_2$}}%
  \end{picture}
\end{equation}
这里四个箭头 $P_1,\cdots,P_4$ 相应于(15.7)式最后一项中的四个 $G$ 因子，而图的“内”部定义了(第一级的)顶角函数——(15.9)图等式左边的小圆.将这些图展成解析形式，得：
\begin{equation*}
  \Gamma^{(1)}_{\gamma\delta,\alpha\beta}(P_3,P_4;P_1,P_2)
  =-\delta_{\alpha\gamma}\delta_{\beta\delta}U(P_1-P_3)
  +\delta_{\alpha\delta}\delta_{\beta\gamma}U(P_1-P_4).
\end{equation*}
更高级的图包含三种修正：1)对两条未相连的实线的进一步修正，2)对(15.9)各图外线的自能型修正，3)还有一种修正，用来构成子图以代替(15.9)图上的虚线；一切可能的这些子图之和便给出精确的顶角函数 $\mathrm{i}\Gamma$.现在以骨架图之和将双粒子格林函数图示出来：
\begin{equation}
  \setlength{\unitlength}{0.72pt}%
  \begin{picture}(560,175)(0,0)
    \put(110,106){\makebox(0,0){$P_3=P_1$}}%
    \tarrL{160}{82}{80}%
    \put(110,58){\makebox(0,0){$P_4=P_2$}}%
    \put(186,82){\makebox(0,0){$+$}}%
    \put(262,106){\makebox(0,0){$P_4=P_1$}}%
    \tarrL{312}{82}{80}%
    \put(262,58){\makebox(0,0){$P_3=P_2$}}%
    \put(338,82){\makebox(0,0){$+$}}%
    \put(390,100){\vector(-1,1){40}}%
    \put(390,64){\vector(-1,-1){40}}%
    \qbezier(424,58)(445.936,60.064)(448,82)%
    \qbezier(448,82)(445.936,103.936)(424,106)%
    \qbezier(424,106)(402.064,103.936)(400,82)%
    \qbezier(400,82)(402.064,60.064)(424,58)%
    \put(424,82){\makebox(0,0){$\mathrm{i}\Gamma$}}%
    \put(482,118){\vector(-1,-1){33}}%
    \put(482,46){\vector(-1,1){33}}%
    \put(356,162){\makebox(0,0){$P_3$}}%
    \put(496,162){\makebox(0,0){$P_1$}}%
    \put(356,4){\makebox(0,0){$P_4$}}%
    \put(496,4){\makebox(0,0){$P_2$}}%
  \end{picture}
\end{equation}
粗线表示精确的 $G$ 函数，而小圆代表顶角函数.

在各级微扰论中计算顶角函数时，应按照 §13 所说的图技术规则进行，而且应当研究四条外线的图(而不研究在计算函数 $G$ 时的两条外线的图).对决定图的共同符号规则 3)，应当补充下述规定：如果外线 1 与 4，2 与 3(代替 1 与 3，2 与 4)是以实线连续依次连接的，那么图就改为相反的符号.

作为例子，我们把确定二级微扰论中顶角函数的图画出来：
\begin{equation}
  \setlength{\unitlength}{0.57pt}%
  \begin{picture}(780,270)(0,0)
    \put(72,208){\vector(-1,1){38}}%
    \put(72,172){\vector(-1,-1){38}}%
    \qbezier(90,164)(113.764,166.236)(116,190)%
    \qbezier(116,190)(113.764,213.764)(90,216)%
    \qbezier(90,216)(66.236,213.764)(64,190)%
    \qbezier(64,190)(66.236,166.236)(90,164)%
    \put(90,190){\makebox(0,0){$\mathrm{i}\Gamma^{(2)}$}}%
    \put(150,238){\vector(-1,-1){30}}%
    \put(152,140){\vector(-1,1){32}}%
    \put(24,254){\makebox(0,0){$P_3$}}%
    \put(164,252){\makebox(0,0){$P_1$}}%
    \put(24,126){\makebox(0,0){$P_4$}}%
    \put(164,128){\makebox(0,0){$P_2$}}%
    \put(192,190){\makebox(0,0){$=$}}%
    \put(266,222){\vector(-1,1){26}}%
    \put(312,252){\vector(-1,-1){28}}%
    \dvh{266}{162}{8}%
    \put(266,154){\vector(-1,-1){26}}%
    \put(312,124){\vector(-1,1){28}}%
    \put(232,258){\makebox(0,0){$3$}}%
    \put(322,258){\makebox(0,0){$1$}}%
    \put(232,120){\makebox(0,0){$4$}}%
    \put(322,120){\makebox(0,0){$2$}}%
    \put(266,96){\makebox(0,0){(a)}}%
    \put(352,190){\makebox(0,0){$+$}}%
    \put(426,222){\vector(-1,1){26}}%
    \put(472,252){\vector(-1,-1){28}}%
    \put(426,222){\vector(-1,-1){40}}%
    \put(470,180){\vector(-1,1){40}}%
    \multiput(398,196)(7,0){10}{\line(1,0){5}}%
    \put(388,258){\makebox(0,0){$3$}}%
    \put(482,258){\makebox(0,0){$1$}}%
    \put(376,168){\makebox(0,0){$4$}}%
    \put(482,166){\makebox(0,0){$2$}}%
    \put(428,96){\makebox(0,0){(b)}}%
    \put(510,190){\makebox(0,0){$+$}}%
    \put(586,222){\vector(-1,1){26}}%
    \put(632,252){\vector(-1,-1){28}}%
    \put(586,222){\line(0,-1){12}}%
    \qbezier(586,194)(593.312,194.688)(594,202)%
    \qbezier(594,202)(593.312,209.312)(586,210)%
    \qbezier(586,210)(578.688,209.312)(578,202)%
    \qbezier(578,202)(578.688,194.688)(586,194)%
    \dvh{586}{140}{7}%
    \put(586,132){\vector(-1,-1){26}}%
    \put(632,102){\vector(-1,1){28}}%
    \put(548,258){\makebox(0,0){$3$}}%
    \put(642,258){\makebox(0,0){$1$}}%
    \put(548,98){\makebox(0,0){$4$}}%
    \put(642,98){\makebox(0,0){$2$}}%
    \put(586,72){\makebox(0,0){(c)}}%
    \put(28,60){\makebox(0,0){$+$}}%
    \put(150,130){\vector(-1,1){26}}%
    \put(196,160){\vector(-1,-1){28}}%
    \put(150,130){\line(-1,-1){22}}%
    \put(150,130){\line(1,-1){22}}%
    \dvh{128}{64}{7}%
    \dvh{172}{64}{7}%
    \arrL{232}{60}{58}%
    \arrL{280}{60}{44}%
    \put(112,166){\makebox(0,0){$3$}}%
    \put(204,166){\makebox(0,0){$1$}}%
    \put(100,40){\makebox(0,0){$4$}}%
    \put(204,40){\makebox(0,0){$2$}}%
    \put(150,18){\makebox(0,0){(d)}}%
    \put(312,60){\makebox(0,0){$+$}}%
    \put(350,114){\line(1,-1){84}}%
    \put(350,46){\line(1,1){84}}%
    \put(350,114){\vector(-1,1){16}}%
    \put(450,30){\vector(-1,1){16}}%
    \put(450,130){\vector(-1,-1){16}}%
    \put(350,46){\vector(-1,-1){16}}%
    \qbezier(392,69)(402.054,69.946)(403,80)%
    \qbezier(403,80)(402.054,90.054)(392,91)%
    \qbezier(392,91)(381.946,90.054)(381,80)%
    \qbezier(381,80)(381.946,69.946)(392,69)%
    \multiput(388,72)(7,0){2}{\line(1,1){5}}%
    \put(326,140){\makebox(0,0){$3$}}%
    \put(458,140){\makebox(0,0){$1$}}%
    \put(326,26){\makebox(0,0){$4$}}%
    \put(458,26){\makebox(0,0){$2$}}%
    \put(392,4){\makebox(0,0){(e)}}%
    \put(490,80){\makebox(0,0)[l]{$+(3\leftrightarrow4)$}}%
  \end{picture}
\end{equation}
自能函数与顶角函数($\Sigma$ 与 $\Gamma$)不是不相关的；它们之间以一定的积分方程(所谓戴森方程)\footnote{该方程类似于量子电动力学中的戴森(Dyson)方程(见第四卷 §107).}相联系.

为推导戴森方程，我们可利用方程(9.5)，该方程(该处曾指出)在考虑到粒子相互作用时也成立.但是与 §9 中的推导相比，其差别在于现在 $\psi$ 算符满足方程(7.8).在此方程中略去含外场的项，并将其中的微商 $\partial\hat{\Psi}/\partial t_1$ 代入(9.5)式，得：
\begin{equation}
\begin{aligned}
  &\Bigl(\mathrm{i}\frac{\partial}{\partial t_1}+\frac{\Delta_1}{2m}+\mu\Bigr)G_{\alpha\beta}(X_1-X_2)
  -\delta_{\alpha\beta}\delta^{(4)}(X_1-X_2)=\\
  =&\ \mathrm{i}\int\langle T\hat{\bm{\Psi}}^{+}_\gamma(X_3)U(X_1-X_3)\hat{\bm{\Psi}}_\gamma(X_3)\mathrm{d}^4X_3\cdot\hat{\bm{\Psi}}_\alpha(X_1)\hat{\bm{\Psi}}^{+}_\beta(X_2)\rangle=\\
  =&-\mathrm{i}\int K_{\gamma\alpha,\gamma\beta}(X_3,X_1;X_3,X_2)U_{13}\,\mathrm{d}^4X_3.
\end{aligned}
\end{equation}
因为 $K$ 可以按照(15.7)式通过 $\Gamma$ 表达出来，所以这个等式原则上解决了所提出的问题.下面要做的只是再把它转到动量表象.为此，给等式(15.12)乘以 $\exp[\mathrm{i}P(X_1-X_2)]$，把 $K_{31,32}$ 和 $U_{13}$ 分别写成(15.5)和(13.9)的形式，并对 $\mathrm{d}^4(X_1-X_2)$ 求积分.这样，对 4 维坐标求积分给出 $\delta$ 函数，后者将因对 4 维动量的积分而消掉.最后得到：
\begin{equation}
\begin{aligned}
  &\bigl[G^{(0)}{}^{-1}(P)G(P)-1\bigr]\delta_{\alpha\beta}=\\
  =&-\mathrm{i}\int K_{\gamma\alpha,\gamma\beta}(P_3,P_4;P_3+P_4-P,P)U(P-P_4)
  \frac{\mathrm{d}^4P_3\mathrm{d}^4P_4}{(2\pi)^8},
\end{aligned}
\end{equation}
式中 $G^{(0)}(P)$ 由(9.7)式得出.

现在还要用 $\Gamma$ 表达 $K$.将(15.7)式代入(15.13)式，最后便得出下列形式的戴森方程：
\begin{equation}
\begin{aligned}
  \delta_{\alpha\beta}\bigl[G^{(0)}{}^{-1}(P)-G^{-1}(P)\bigr]&=\delta_{\alpha\beta}\Sigma(P)=\\
  ={}&U(0)n(\mu)\delta_{\alpha\beta}+\mathrm{i}\delta_{\alpha\beta}\int U(P-P_1)G(P_1)\frac{\mathrm{d}^4P_1}{(2\pi)^4}+\\
  +{}&\int\Gamma_{\gamma\alpha,\gamma\beta}(P_3,P_4;P_3+P_4-P,P)G(P_3)G(P_4)G(P_3+P_4-P)\times\\
  \times{}&U(P-P_4)\frac{\mathrm{d}^4P_3\mathrm{d}^4P_4}{(2\pi)^8},
\end{aligned}
\end{equation}
这里 $n(\mu)$ 为系统的精确密度，它是系统化学势的函数；这个因子是按公式(7.24)积分 $G$ 函数时出现的(同时考虑该 $G$ 函数是由收缩 $\hat{\Psi}^{+}\hat{\Psi}$ 产生的).注意，方程(15.14)右边的第一项就是(14.11)式的 $\Sigma_a$.

\section{顶角函数与准粒子散射幅的关系}

以上几节建立的数学工具，使我们有可能严格论证和更深刻理解朗道费米液体理论基本关系式的意义，这些关系式在第一章里曾用直观方法在一定程度上作过介绍.在 §16— §20 就将论述这一问题\footnote{§16—§18 的内容是属于 Л. Д. 朗道的工作(1958)，§19，§20 的内容是属于 Л. Д. 朗道和 Л. П. 皮塔耶夫斯基的工作(1959).}.

顶角函数与准粒子相互散射幅之间存在密切的关系.为了更好地了解这种关系，我们首先在纯量子力学范围内研究真空中两个粒子的散射问题.

在量子力学中，“四条腿图”是有四条外线(两条入射线和两条出射线)的图——相应于两个粒子的碰撞过程；而在图的解析表达式中，外线相当于自由粒子波函数(平面波)的振幅(见第四卷 §106).现在，我们将看到各级图怎样如实地给出散射幅通常的非相对论玻恩展开序列的各项.

首先，在真空情况下大多数图总是化为零的.这一点，在坐标表象中最容易理解，因为我们注意到：在真空中一切 $\langle\hat{\Psi}^{+}\hat{\Psi}\rangle$ 形的收缩都等于零，此时湮没算符位于右边并首先作用在真空态上，而剩余的只是 $\langle\hat{\Psi}\hat{\Psi}^{+}\rangle$ 形的收缩项.因此，一切含实线圈的图都变成零，这些图总有 $\langle\hat{\Psi}^{+}\hat{\Psi}\rangle$ 形的收缩.根据同样理由，对格林函数(也就是对图的实内线)的一切修正都等于零\footnote{在真空中格林函数不出现任何修正只表明单个粒子跟任何别的都不发生相互作用. 与此相关，我们记得相对论粒子格林函数的真空修正与可能出现虚电子对和虚光子的中间态有关.}.最后，具有交叉虚线的图也为零；例如，在下图中(这里数字 1 和 2 表示宗量 $t_1$ 和 $t_2$)，当 $t_2>t_1$ 时，上内线对应于收缩 $\langle\hat{\Psi}^{+}_2\hat{\Psi}_1\rangle=0$，如果 $t_2<t_1$，则收缩 $\langle\hat{\Psi}^{+}_1\hat{\Psi}_2\rangle=0$ 对应于下内线.
\begin{equation*}
  \setlength{\unitlength}{0.80pt}%
  \begin{picture}(240,100)(0,0)
    \arrL{200}{72}{80}%
    \arrL{200}{28}{80}%
    \put(122,70){\line(1,-1){10}}%
    \put(136,54){\line(1,-1){10}}%
    \put(150,40){\line(1,-1){8}}%
    \put(178,70){\line(-1,-1){10}}%
    \put(164,54){\line(-1,-1){10}}%
    \put(150,40){\line(-1,-1){8}}%
    \put(106,84){\makebox(0,0){$1$}}%
    \put(194,84){\makebox(0,0){$2$}}%
    \put(106,16){\makebox(0,0){$2$}}%
    \put(194,16){\makebox(0,0){$1$}}%
  \end{picture}
\end{equation*}
这样一来，对于真空中的两个粒子，只剩下如下的图，它们组成所谓“梯形级数”：
\begin{equation}
  \setlength{\unitlength}{0.70pt}%
  \begin{picture}(640,116)(0,0)
    \put(58,72){\vector(-1,1){28}}%
    \put(58,48){\vector(-1,-1){28}}%
    \qbezier(70,42)(86.452,43.548)(88,60)%
    \qbezier(88,60)(86.452,76.452)(70,78)%
    \qbezier(70,78)(53.548,76.452)(52,60)%
    \qbezier(52,60)(53.548,43.548)(70,42)%
    \put(70,60){\makebox(0,0){$\mathrm{i}\Gamma$}}%
    \put(116,98){\vector(-1,-1){28}}%
    \put(116,22){\vector(-1,1){28}}%
    \put(20,108){\makebox(0,0){$P_3$}}%
    \put(126,108){\makebox(0,0){$P_1$}}%
    \put(20,10){\makebox(0,0){$P_4$}}%
    \put(126,10){\makebox(0,0){$P_2$}}%
    \put(146,60){\makebox(0,0){$=$}}%
    \arrL{266}{84}{60}%
    \arrL{266}{36}{60}%
    \dvh{236}{42}{6}%
    \put(286,60){\makebox(0,0){$+$}}%
    \arrL{386}{84}{60}%
    \arrL{386}{36}{60}%
    \dvh{340}{42}{6}%
    \dvh{370}{42}{6}%
    \put(406,60){\makebox(0,0){$+$}}%
    \arrL{506}{84}{60}%
    \arrL{506}{36}{60}%
    \dvh{460}{42}{6}%
    \dvh{480}{42}{6}%
    \dvh{500}{42}{6}%
    \put(526,60){\makebox(0,0)[l]{$+\cdots+(3\leftrightarrow4)$}}%
  \end{picture}
\end{equation}
图中的实内线对应于真空格林函数
\begin{equation}
  G^{(\text{真空})}(\omega,\bm{p})=\Bigl[\omega-\frac{p^2}{2m}+\mathrm{i}0\Bigr]^{-1}
\end{equation}
［这正是 $\mu=0$ 时的(9.7)式］.应注意，由于分母没有 $\mu$，这个函数的极点总在复变数 $\omega$ 的一定的(下)半平面上.上述图都变为零的情况之所以发生，从数学观点来说，是由于被积式的全部极点都分布在同一个半平面上的缘故.在另一半平面上闭合积分路径时，积分显然变为零.

将梯形级数(16.1)转化为积分方程，便可以对该级数求和［对照下面类似的级数(17.3)求和］.如果首先略去具有交换外线端点 3 和 4 的图，则该方程与动量表象中不考虑全同性的双粒子薛定谔方程［第三卷方程(130.9)］是等价的.相应地，顶角函数 $\Gamma$ 可用下式通过双粒子的散射幅 $f$ 表达出来：
\begin{equation}
  \Gamma_{\gamma\delta,\alpha\beta}(P_3,P_4;P_1,P_2)=\delta_{\alpha\gamma}\delta_{\beta\delta}\frac{4\pi}{m}f.
\end{equation}
若添上具有交换外线端点 3 和 4 的图，会使散射幅反对称化，对于费米子本应这样.在微扰论的一级近似下，只剩下(16.1)的第一图和具有交换外线端点的一个图，其中完全不含 $G^{(\text{真空})}$.于是可得到散射幅通常的一级玻恩近似公式.以后的高级图，在对中间的频率进行积分之后，将给出散射幅修正的高级玻恩近似的已知表达式.

在费米液体中，与介质粒子相互作用的碰撞粒子可等效地用准粒子来代替.所有与这种相互作用有关的对图的内线修正，函数 $\Gamma$ 的定义都自动地顾及到了.但是对外线修正，需另加考虑.量子场论指出，由于普遍要求散射矩阵具有么正性，这些修正将使散射幅中的每条自由外线出现因子 $\sqrt{Z}$，这里 $Z$ 是格林函数的重整化常数(见第四卷 §110)；对于具有四条外线的图，这意味着要乘以 $Z^2$.虽说量子场论的结论对于费米液体中的准粒子也成立，但在这里我们也可以借助更简单(虽然并不严格)的讨论来解释这个乘数的来由.

问题在于，液体的格林函数在接近自己的极点处［(10.2)式中的第一项］与理想气体的格林函数只差一个因子 $Z$.如果以算符 $\hat{\Psi}_{\mathrm{qu}}=\hat{\Psi}/\sqrt{Z}$，$\hat{\Psi}^{+}_{\mathrm{qu}}=\hat{\Psi}^{+}/\sqrt{Z}$ 分别代替 $\hat{\Psi}$ 和 $\hat{\Psi}^{+}$，那么由它们组成的格林函数 $G_{\mathrm{qu}}=G/Z$，就很像理想气体的格林函数在极点附近的行为.就此意义来说，这些算符可以看作是准粒子理想气体的 $\psi$ 算符.按这些算符确定的双粒子格林函数将是 $K_{\mathrm{qu}}=K/Z^2$，因此［依照定义(15.7)］顶角部分 $\Gamma_{\mathrm{qu}}=\Gamma Z^2$，这就是我们所需要的证明.

在应用于准粒子时，我们感兴趣的与其说是散射截面，不如说是(1\,s 内在 $1\,\mathrm{cm}^3$ 液体中的)碰撞数.对于粒子的动量变化和自旋投影变化($\bm{p}_1\alpha$，$\bm{p}_2\beta\to\bm{p}_3\gamma$，$\bm{p}_4\delta$)都已知的碰撞来说，这种碰撞数用下列公式给出：
\begin{equation}
\begin{aligned}
  \mathrm{d}W={}&2\pi\bigl|Z^2\Gamma_{\gamma\delta,\alpha\beta}(P_3,P_4;P_1,P_2)\bigr|^2\delta(\varepsilon_3+\varepsilon_4-\varepsilon_1-\varepsilon_2)\times\\
  &\qquad\times n_{\bm{p}_1}n_{\bm{p}_2}(1-n_{\bm{p}_3})(1-n_{\bm{p}_4})\frac{\mathrm{d}^3p_1\mathrm{d}^3p_2\mathrm{d}^3p_3}{(2\pi)^9},
\end{aligned}
\end{equation}
其中 $\bm{p}_1+\bm{p}_2=\bm{p}_3+\bm{p}_4$，$n_{\bm{p}}$ 是准粒子的分布函数.$n_{\bm{p}_1}$ 和 $n_{\bm{p}_2}$ 两个因子只是表明这一事实：初动量(和自旋投影)为已知的准粒子的碰撞数与单位体积中这种准粒子数成正比.根据泡利原理，因子 $(1-n_{\bm{p}_3})$ 和 $(1-n_{\bm{p}_4})$ 与下列事实有关：只有末态未被占据时，碰撞才能发生.

\section{小动量传递时的顶角函数}

在费米液体理论中，当两对变量 $P_1$ 与 $P_3$ 及 $P_2$ 与 $P_4$ 的值都很接近时，顶角函数起着重要作用(其中我们将看到该函数与准粒子相互作用函数有密切关系).鉴于下列关系：
\begin{equation*}
  P_1+P_2=P_3+P_4,
\end{equation*}
假定 $P_3=P_1+K$，$P_4=P_2-K$，并引入简化记号
\begin{equation}
  \Gamma_{\gamma\delta,\alpha\beta}(P_1+K,P_2-K;P_1,P_2)=\Gamma_{\gamma\delta,\alpha\beta}(K;P_1,P_2),
\end{equation}
我们将研究这个函数在小 $K$ 值的情况.用准粒子散射过程的术语来说，就是要研究接近于“向前散射”的碰撞如何传递小的 4 维动量.

当 $K=0$ 时，我们看到函数 $\Gamma$ 具有奇异性；而我们感兴趣的，正是该函数具有这种奇异性的那一部分.根据对如下骨架图的分析，不难理解这种奇异性的来由：
\begin{equation}
  \setlength{\unitlength}{0.85pt}%
  \begin{picture}(310,110)(0,0)
    \put(58,88){\vector(1,-1){26}}%
    \put(232,88){\vector(-1,-1){26}}%
    \qbezier(90,46)(98.226,46.774)(99,55)%
    \qbezier(99,55)(98.226,63.226)(90,64)%
    \qbezier(90,64)(81.774,63.226)(81,55)%
    \qbezier(81,55)(81.774,46.774)(90,46)%
    \qbezier(200,46)(208.226,46.774)(209,55)%
    \qbezier(209,55)(208.226,63.226)(200,64)%
    \qbezier(200,64)(191.774,63.226)(191,55)%
    \qbezier(191,55)(191.774,46.774)(200,46)%
    \qbezier(90,64)(145,88)(200,64)%
    \put(145,96){\makebox(0,0){$Q{+}K$}}%
    \tarrL{200}{46}{70}%
    \put(165,30){\makebox(0,0){$Q$}}%
    \put(84,48){\vector(-1,-1){26}}%
    \put(206,48){\vector(1,-1){26}}%
    \put(46,96){\makebox(0,0){$P_1$}}%
    \put(244,96){\makebox(0,0){$P_2$}}%
    \put(42,12){\makebox(0,0){$P_1{+}K$}}%
    \put(250,12){\makebox(0,0){$P_2{-}K$}}%
  \end{picture}
\end{equation}
骨架图包含双粒子格林函数一些图的集合，这些图都能在两对外点 $P_1,P_3$ 和 $P_2,P_4$ 之间被截成两段，它们原用两条实线相连\footnote{例如，在(17.2)式按成对相互作用的二级微扰论中，包含图(15.11a,b,c)及交换外线端点 3 和 4 的图(15.11e).}.两条连结的粗线相应于精确的单粒子格林函数 $G(Q)$ 和 $G(Q+K)$，并且要对图中 4 维动量 $Q$ 进行积分.当 $K\to0$ 时，这两个函数的宗量相互靠近，因而它们的极点也相互靠近.相互靠近的极点将“紧夹”积分路线(见下文)，这就是函数 $\Gamma$ 中出现奇异性的根源.

为了算出精确的函数 $\Gamma$，需要对微扰论的整个级数求和.由于我们的目的在于分离出 $K=0$ 时具有奇异性的部分，首先要将下述所有图的贡献分离开，即这些图不能靠切断两条 4 维动量相接近(只差 $K$ 值)的实线而被截开.当 $K=0$ 时函数 $\Gamma$ 不具有奇异性的部分，我们记作 $\tilde{\Gamma}$；在 $\tilde{\Gamma}$ 中可以设 $K=0$，因此 $\tilde{\Gamma}$ 只是变量 $P_1,P_2$ 的函数：$\tilde{\Gamma}_{\gamma\delta,\alpha\beta}(P_1,P_2)$.至于说到“危险”的图，可将它们按其所含宗量相近的双线对的数目来分类.因此，完全的顶角部分 $\Gamma$ 可用下图的无穷“梯形”级数表达：
\begin{equation}
  \setlength{\unitlength}{0.70pt}%
  \begin{picture}(620,115)(0,0)
    \put(34,92){\vector(1,-1){30}}%
    \put(106,92){\vector(-1,-1){30}}%
    \put(70,60){\vector(-1,-1){30}}%
    \put(72,58){\vector(1,-1){30}}%
    \put(24,100){\makebox(0,0){$P_1$}}%
    \put(118,100){\makebox(0,0){$P_2$}}%
    \put(22,18){\makebox(0,0){$P_1{+}K$}}%
    \put(122,18){\makebox(0,0){$P_2{-}K$}}%
    \put(142,60){\makebox(0,0){$=$}}%
    \put(158,94){\vector(1,-1){26}}%
    \put(232,94){\vector(-1,-1){26}}%
    \put(187,53){\vector(-1,-1){30}}%
    \put(203,53){\vector(1,-1){30}}%
    \qbezier(195,49)(205.054,49.946)(206,60)%
    \qbezier(206,60)(205.054,70.054)(195,71)%
    \qbezier(195,71)(184.946,70.054)(184,60)%
    \qbezier(184,60)(184.946,49.946)(195,49)%
    \put(189,53){\line(1,1){12}}%
    \put(189,60){\line(1,1){7}}%
    \put(148,102){\makebox(0,0){$P_1$}}%
    \put(242,102){\makebox(0,0){$P_2$}}%
    \put(144,10){\makebox(0,0){$P_1{+}K$}}%
    \put(246,10){\makebox(0,0){$P_2{-}K$}}%
    \put(256,60){\makebox(0,0){$+$}}%
    \qbezier(300,49)(310.054,49.946)(311,60)%
    \qbezier(311,60)(310.054,70.054)(300,71)%
    \qbezier(300,71)(289.946,70.054)(289,60)%
    \qbezier(289,60)(289.946,49.946)(300,49)%
    \qbezier(368,49)(378.054,49.946)(379,60)%
    \qbezier(379,60)(378.054,70.054)(368,71)%
    \qbezier(368,71)(357.946,70.054)(357,60)%
    \qbezier(357,60)(357.946,49.946)(368,49)%
    \put(294,54){\line(1,1){11}}%
    \put(294,62){\line(1,1){6}}%
    \put(362,54){\line(1,1){11}}%
    \put(362,62){\line(1,1){6}}%
    {\thicklines\qbezier(300,71)(334,90)(368,71)}%
    {\thicklines\qbezier(300,49)(334,30)(368,49)}%
    \put(328,86){\vector(1,0){12}}%
    \put(340,34){\vector(-1,0){12}}%
    \put(398,60){\makebox(0,0){$+$}}%
    \qbezier(440,50)(449.14,50.86)(450,60)%
    \qbezier(450,60)(449.14,69.14)(440,70)%
    \qbezier(440,70)(430.86,69.14)(430,60)%
    \qbezier(430,60)(430.86,50.86)(440,50)%
    \qbezier(485,50)(494.14,50.86)(495,60)%
    \qbezier(495,60)(494.14,69.14)(485,70)%
    \qbezier(485,70)(475.86,69.14)(475,60)%
    \qbezier(475,60)(475.86,50.86)(485,50)%
    \qbezier(530,50)(539.14,50.86)(540,60)%
    \qbezier(540,60)(539.14,69.14)(530,70)%
    \qbezier(530,70)(520.86,69.14)(520,60)%
    \qbezier(520,60)(520.86,50.86)(530,50)%
    \put(435,55){\line(1,1){10}}%
    \put(435,62){\line(1,1){5}}%
    \put(480,55){\line(1,1){10}}%
    \put(480,62){\line(1,1){5}}%
    \put(525,55){\line(1,1){10}}%
    \put(525,62){\line(1,1){5}}%
    {\thicklines\qbezier(440,70)(462,86)(485,70)}%
    {\thicklines\qbezier(440,50)(462,34)(485,50)}%
    {\thicklines\qbezier(485,70)(507,86)(530,70)}%
    {\thicklines\qbezier(485,50)(507,34)(530,50)}%
    \put(458,82){\vector(1,0){10}}%
    \put(512,38){\vector(-1,0){10}}%
    \put(548,60){\makebox(0,0)[l]{$+\cdots$}}%
  \end{picture}
\end{equation}
这里，空白小圆相应于待求的 $\mathrm{i}\Gamma$，阴影小圆代表 $\mathrm{i}\tilde{\Gamma}$.这些图上的外线，不含在 $\tilde{\Gamma}$ 的定义中，它们只用于指明入射和出射的 4 维动量的数目和数值.

图(17.3)上的一切内线都是粗线，它们对应于精确的 $G$ 函数.因此我们强调，之所以可将 $\Gamma$ 表现为这些骨架图(以及由之得出的一切推论)，绝不是假定粒子之间具有成对相互作用，因为这里没有显形式的虚线，实际上只有小圆所表示的单元内部构造(这里我们不感兴趣)才与相互作用性质有关\footnote{这里只假定有如粒子数守恒一类的普遍性质. 粒子数守恒表现在图的每个断面上向右通过和向左通过的线数之差保持不变(对于(17.3)型图的各个断面，这个差等于零).}.

对级数(17.3)的求和问题，归结为解积分方程，为了得到这个方程，整个级数尚需“乘”一个 $\tilde{\Gamma}$，即以如下级数代替(17.3)：
\begin{equation*}
  \setlength{\unitlength}{0.70pt}%
  \begin{picture}(620,110)(0,0)
    \put(28,90){\vector(1,-1){26}}%
    \put(94,90){\vector(-1,-1){26}}%
    \put(60,60){\vector(-1,-1){26}}%
    \put(62,58){\vector(1,-1){26}}%
    \qbezier(110,50)(119.14,50.86)(120,60)%
    \qbezier(120,60)(119.14,69.14)(110,70)%
    \qbezier(110,70)(100.86,69.14)(100,60)%
    \qbezier(100,60)(100.86,50.86)(110,50)%
    \put(106,55){\line(1,1){10}}%
    \qbezier(175,50)(184.14,50.86)(185,60)%
    \qbezier(185,60)(184.14,69.14)(175,70)%
    \qbezier(175,70)(165.86,69.14)(165,60)%
    \qbezier(165,60)(165.86,50.86)(175,50)%
    \qbezier(110,70)(142,88)(175,70)%
    \qbezier(110,50)(142,32)(175,50)%
    \put(220,60){\makebox(0,0){$=$}}%
    \qbezier(260,50)(269.14,50.86)(270,60)%
    \qbezier(270,60)(269.14,69.14)(260,70)%
    \qbezier(260,70)(250.86,69.14)(250,60)%
    \qbezier(250,60)(250.86,50.86)(260,50)%
    \qbezier(325,50)(334.14,50.86)(335,60)%
    \qbezier(335,60)(334.14,69.14)(325,70)%
    \qbezier(325,70)(315.86,69.14)(315,60)%
    \qbezier(315,60)(315.86,50.86)(325,50)%
    \put(256,55){\line(1,1){10}}%
    \put(321,55){\line(1,1){10}}%
    \qbezier(260,70)(292,88)(325,70)%
    \qbezier(260,50)(292,32)(325,50)%
    \put(352,60){\makebox(0,0){$+$}}%
    \qbezier(395,50)(404.14,50.86)(405,60)%
    \qbezier(405,60)(404.14,69.14)(395,70)%
    \qbezier(395,70)(385.86,69.14)(385,60)%
    \qbezier(385,60)(385.86,50.86)(395,50)%
    \qbezier(440,50)(449.14,50.86)(450,60)%
    \qbezier(450,60)(449.14,69.14)(440,70)%
    \qbezier(440,70)(430.86,69.14)(430,60)%
    \qbezier(430,60)(430.86,50.86)(440,50)%
    \qbezier(485,50)(494.14,50.86)(495,60)%
    \qbezier(495,60)(494.14,69.14)(485,70)%
    \qbezier(485,70)(475.86,69.14)(475,60)%
    \qbezier(475,60)(475.86,50.86)(485,50)%
    \put(391,55){\line(1,1){10}}%
    \put(436,55){\line(1,1){10}}%
    \put(481,55){\line(1,1){10}}%
    \qbezier(395,70)(417,86)(440,70)%
    \qbezier(395,50)(417,34)(440,50)%
    \qbezier(440,70)(462,86)(485,70)%
    \qbezier(440,50)(462,34)(485,50)%
    \put(510,60){\makebox(0,0)[l]{$+\cdots$}}%
  \end{picture}
\end{equation*}
与原来的级数(17.3)比较，可得到等式
\begin{equation}
  \setlength{\unitlength}{0.72pt}%
  \begin{picture}(640,120)(0,0)
    \put(58,88){\vector(1,-1){26}}%
    \put(232,88){\vector(-1,-1){26}}%
    \qbezier(95,50)(117,60)(95,70)%
    \qbezier(95,50)(73,60)(95,70)%
    \put(91,55){\line(1,1){10}}%
    \qbezier(200,50)(222,60)(200,70)%
    \qbezier(200,50)(178,60)(200,70)%
    \qbezier(95,70)(148,94)(200,70)%
    \qbezier(95,50)(148,26)(200,50)%
    \put(148,102){\makebox(0,0){$Q{+}K$}}%
    \put(148,18){\makebox(0,0){$Q$}}%
    \put(46,96){\makebox(0,0){$P_1$}}%
    \put(244,96){\makebox(0,0){$P_2$}}%
    \put(84,54){\vector(-1,-1){26}}%
    \put(206,54){\vector(1,-1){26}}%
    \put(34,12){\makebox(0,0){$P_1{+}K$}}%
    \put(256,12){\makebox(0,0){$P_2{-}K$}}%
    \put(286,60){\makebox(0,0){$=$}}%
    \put(328,88){\vector(1,-1){26}}%
    \put(372,88){\vector(-1,-1){26}}%
    \put(350,60){\vector(-1,-1){26}}%
    \put(352,58){\vector(1,-1){26}}%
    \put(318,96){\makebox(0,0){$P_1$}}%
    \put(384,96){\makebox(0,0){$P_2$}}%
    \put(316,12){\makebox(0,0){$P_1{+}K$}}%
    \put(386,12){\makebox(0,0){$P_2{-}K$}}%
    \put(420,60){\makebox(0,0){$-$}}%
    \qbezier(495,50)(517,60)(495,70)%
    \qbezier(495,50)(473,60)(495,70)%
    \put(491,55){\line(1,1){10}}%
    \put(491,62){\line(1,1){5}}%
    \put(460,88){\vector(1,-1){26}}%
    \put(534,88){\vector(-1,-1){26}}%
    \put(489,54){\vector(-1,-1){30}}%
    \put(503,54){\vector(1,-1){30}}%
    \put(450,96){\makebox(0,0){$P_1$}}%
    \put(546,96){\makebox(0,0){$P_2$}}%
    \put(448,12){\makebox(0,0){$P_1{+}K$}}%
    \put(548,12){\makebox(0,0){$P_2{-}K$}}%
  \end{picture}
\end{equation}
将这个图等式写成解析形式时，即给出待求的积分方程：
\begin{equation}
\begin{aligned}
  \Gamma_{\gamma\delta,\alpha\beta}(K;&P_1,P_2)=\\
  ={}&\tilde{\Gamma}_{\gamma\delta,\alpha\beta}(P_1,P_2)-\mathrm{i}\int\tilde{\Gamma}_{\gamma\zeta,\alpha\kappa}(P_1,Q)G(Q+K)G(Q)\times\\
  &\times\Gamma_{\kappa\delta,\zeta\beta}(K;Q,P_2)\frac{\mathrm{d}^4Q}{(2\pi)^4}.
\end{aligned}
\end{equation}
这里，如上所述，在函数 $\tilde{\Gamma}$ 中假定 $K=0$；利用前面已引入的 $\Gamma$ 和 $\tilde{\Gamma}$ 的简化记号，并假定 $G_{\alpha\beta}=G\delta_{\alpha\beta}$.

为了考察这个方程，我们首先研究方程核中的乘积 $G(Q+K)G(Q)$.我们讲过，当 $K$ 小时，两个因子的极点相互接近.这些极点附近 $G$ 函数表现为极点项(10.2).把 4 维矢量 $K$ 和 $Q$ 的分量按照
\begin{equation}
  K=(\omega,\bm{k}),\qquad Q=(q_0,\bm{q})
\end{equation}
表示出来，于是在这个范围内，可写成：
\begin{equation}
\begin{aligned}
  G(Q)G(Q+K)\approx Z^2&\bigl[q_0-v_{\mathrm{F}}(q-p_{\mathrm{F}})+\\
  &+\mathrm{i}\delta_1\bigr]^{-1}\bigl[q_0+\omega-v_{\mathrm{F}}\bigl(|\bm{q}+\bm{k}|-p_{\mathrm{F}}\bigr)+\mathrm{i}\delta_2\bigr]^{-1}
\end{aligned}
\end{equation}
式中 $\delta_1,\delta_2$ 为无穷小附加项，它们的符号(在极点附近)根据下式确定：
\begin{equation}
\begin{aligned}
  \operatorname{sign}\delta_1&=\operatorname{sign}(q-p_{\mathrm{F}}),\\
  \operatorname{sign}\delta_2&=\operatorname{sign}\bigl(|\bm{q}+\bm{k}|-p_{\mathrm{F}}\bigr).
\end{aligned}
\end{equation}
$\delta_1$ 和 $\delta_2$ 的符号决定极点在复变量 $q_0$ 的上半平面或下半平面上的位置.在极点间夹紧 $\mathrm{d}q_0$ 积分的回路(实轴).结果，积分方程的核(随之方程的解)将出现奇异性.极点应位于回路的两侧，即位于不同的半平面上.

我们首先假定 $\bm{q}\cdot\bm{k}>0$，即 $\cos\theta>0$，这里 $\theta$ 为 $\bm q$ 和 $\bm k$ 的夹角.如果 $q<p_{\mathrm{F}}$，$|\bm{q}+\bm{k}|>p_{\mathrm{F}}$，于是 $|\bm{q}+\bm{k}|>q$，并且 $\delta_1$ 和 $\delta_2$ 具有不同的符号($\delta_1<0,\delta_2>0$)，这是由于 $k$ 的微小性等价于
\begin{equation}
  p_{\mathrm{F}}-k\cos\theta<q<p_{\mathrm{F}}.
\end{equation}
在(17.5)式中进一步对 $\mathrm{d}q_0$ 求积分时，可以用无穷远的半圆周(在上半平面或在下半平面全一样)将积分路线封闭起来，这时，积分便决定于被积式在相应的极点的留数.同时，由于(17.9)的间隔狭窄(当 $k$ 微小时)，在积分号下 $\Gamma$ 和 $\tilde{\Gamma}$ 的因子中可取 $k=0$，相应地对于各极点的位置来说(当 $k,\omega$ 微小时)：$q_0\approx0$.

换言之，在积分方程(17.5)核中极点因子的乘积(17.7)，就其本身所起作用的意义来说，相当于 $\delta$ 函数：
\begin{equation*}
  A\delta(q_0)\delta(q-p_{\mathrm{F}}),
\end{equation*}
其中系数 $A$ 由如下积分决定：
\begin{equation*}
  A=\int\frac{Z^2\,\mathrm{d}q_0\,\mathrm{d}q}{\bigl[q_0-v_{\mathrm{F}}(q-p_{\mathrm{F}})+\mathrm{i}\delta_1\bigr]\bigl[q_0+\omega-v_{\mathrm{F}}\bigl(|\bm{q}+\bm{k}|-p_{\mathrm{F}}\bigr)-\mathrm{i}\delta_2\bigr]}.
\end{equation*}
当 $q$ 位于(17.9)间隔之外时，则两个极点位于复变量 $q_0$ 的同一个半平面上，而对 $\mathrm{d}q_0$ 的积分围道在另一个半平面上封闭，我们知道积分将为零.但在(17.9)的区域内，在其中一个半平面上封闭积分围道，并对这个半平面上极点的留数计算积分，我们求得：
\begin{equation*}
  A=\int\frac{2\pi\mathrm{i}Z^2\,\mathrm{d}q}{\omega-v_{\mathrm{F}}\bigl(|\bm{q}+\bm{k}|-q\bigr)+\mathrm{i}0}
\end{equation*}
［注意，在(17.9)的区域内 $\delta_1<0,\delta_2>0$］.因为按条件(17.9)，$q\approx p_{\mathrm{F}}\gg k$，所以可取 $|\bm{q}+\bm{k}|-q\approx k\cos\theta$，再考虑到(17.9)的两侧界限，就有：
\begin{equation*}
  A=\frac{2\pi\mathrm{i}Z^2k\cos\theta}{\omega-kv_{\mathrm{F}}\cos\theta}.
\end{equation*}
用同样方法不难证明，当 $\cos\theta<0$ 时，也可以得到 $A$ 的这个表达式(但 i0 的符号不同，这时应在 $q>p_{\mathrm{F}}$，$|\bm{q}+\bm{k}|<p_{\mathrm{F}}$ 的区域内进行积分).因此，方程(17.5)的核中有：
\begin{equation}
  G(Q)G(Q+K)=\frac{2\pi\mathrm{i}Z^2\,\bm{l}\cdot\bm{k}\,\delta(q_0)\delta\bigl(|\bm{q}|-p_{\mathrm{F}}\bigr)}{\omega-v_{\mathrm{F}}\bm{l}\cdot\bm{k}+\mathrm{i}0\cdot\operatorname{sign}\omega}+\varphi(Q).
\end{equation}
式中写出 $\bm{l}\cdot\bm{k}$ 以代替 $k\cos\theta$($\bm{l}=\bm{q}/q$)，函数 $\varphi$(当小 $K$ 时)不包含 $\delta$ 函数部分，因此其中可以取 $K=0$.

将(17.10)式代入(17.5)式，我们得到下列形式的基本积分方程：
\begin{equation}
\begin{aligned}
  \Gamma_{\gamma\delta,\alpha\beta}(K;P_1,P_2)&=\tilde{\Gamma}_{\gamma\delta,\alpha\beta}(P_1,P_2)-\\
  -\mathrm{i}&\int\tilde{\Gamma}_{\gamma\zeta,\alpha\kappa}(P_1,Q)\varphi(Q)\Gamma_{\kappa\delta,\zeta\beta}(K;Q,P_2)\frac{\mathrm{d}^4Q}{(2\pi)^4}+\\
  +\frac{Z^2p_{\mathrm{F}}^2}{(2\pi)^3}&\int\tilde{\Gamma}_{\gamma\zeta,\alpha\kappa}(P_1,Q_{\mathrm{F}})\Gamma_{\kappa\delta,\zeta\beta}(K;Q_{\mathrm{F}},P_2)\frac{\bm{l}\cdot\bm{k}\,\mathrm{d}o_l}{\omega-v_{\mathrm{F}}\bm{l}\cdot\bm{k}}.
\end{aligned}
\end{equation}
在最后一项中作代换：$\mathrm{d}^4Q=q^2\,\mathrm{d}q\,\mathrm{d}o_l\,\mathrm{d}q_0$(这里 $\mathrm{d}o_l$ 为 $\bm l$ 方向上的立体角元)，并对 $\mathrm{d}q\,\mathrm{d}q_0$ 取积分消去 $\delta$ 函数.在这一项的函数 $\Gamma$ 和 $\tilde{\Gamma}$ 中，宗量 $Q$ 取费米面上的值：$Q_{\mathrm{F}}=(0,p_{\mathrm{F}}\bm{l})$.

应注意到方程(17.11)的核中因子 $\bm{l}\cdot\bm{k}/(\omega-v_{\mathrm{F}}\bm{l}\cdot\bm{k})$ 的特殊性质：当 $\bm{k}\to0$、$\omega\to0$ 时，它的极限与这时比值 $\omega/k$ 所趋近的极限有关.因而方程的解亦将具有这种特性：函数 $\Gamma(K;P_1,P_2)$ 当 $K\to0$ 时的极限依赖于 $\omega$ 和 $k$ 趋近于零的方式.

我们以 $\Gamma^{\omega}(P_1,P_2)$ 表示极限：
\begin{equation}
  \Gamma^{\omega}_{\gamma\delta,\alpha\beta}(P_1,P_2)=\lim_{K\to0}\Gamma_{\gamma\delta,\alpha\beta}(K;P_1,P_2),\ \text{当}\ k/\omega\to0
\end{equation}
(在 §18 中我们将会看到，准粒子相互作用函数恰好与这个量有关).当以这种方式过渡到极限时，(17.11)式中最后一个积分项的核变为零，所以 $\Gamma^{\omega}$ 满足方程：
\begin{equation}
\begin{aligned}
  &\Gamma^{\omega}_{\gamma\delta,\alpha\beta}(P_1,P_2)=\\
  ={}&\tilde{\Gamma}_{\gamma\delta,\alpha\beta}(P_1,P_2)-\mathrm{i}\int\tilde{\Gamma}_{\gamma\zeta,\alpha\kappa}(P_1,Q)\varphi(Q)\Gamma^{\omega}_{\kappa\delta,\zeta\beta}(Q,P_2)\frac{\mathrm{d}^4Q}{(2\pi)^4}.
\end{aligned}
\end{equation}
注意，由于(15.8)式，则
\begin{equation}
  \Gamma^{\omega}_{\gamma\delta,\alpha\beta}(P_1,P_2)=\Gamma^{\omega}_{\delta\gamma,\beta\alpha}(P_2,P_1).
\end{equation}
由(17.11)和(17.13)两个方程可以消掉 $\tilde{\Gamma}$.消掉后，结果为：
\begin{equation}
\begin{aligned}
  \Gamma_{\gamma\delta,\alpha\beta}(K;P_1,P_2)&=\Gamma^{\omega}_{\gamma\delta,\alpha\beta}(P_1,P_2)+\\
  +{}&\frac{Z^2p_{\mathrm{F}}^2}{(2\pi)^3}\int\Gamma^{\omega}_{\gamma\zeta,\alpha\kappa}(P_1,Q_{\mathrm{F}})\Gamma_{\kappa\delta,\zeta\beta}(K;Q_{\mathrm{F}},P_2)\frac{\bm{l}\cdot\bm{k}\,\mathrm{d}o_l}{\omega-v_{\mathrm{F}}\bm{l}\cdot\bm{k}}.
\end{aligned}
\end{equation}
事实上，如果形式上将(17.3)式写成 $\tilde{\Gamma}=\hat{L}\Gamma^{\omega}$，则(17.11)式可写成：
\begin{equation*}
  \hat{L}\Gamma=\tilde{\Gamma}+\frac{Z^2p_{\mathrm{F}}^2}{(2\pi)^3}\int\tilde{\Gamma}\,\Gamma\frac{\bm{l}\cdot\bm{k}}{\omega-v_{\mathrm{F}}\bm{l}\cdot\bm{k}}\,\mathrm{d}o_l
\end{equation*}
将 $\tilde{\Gamma}=\hat{L}\Gamma^{\omega}$ 代入上式，并将算符 $\hat{L}^{-1}$ 作用到等式的两边，便得到(17.15)式.

现在我们按照下式引入函数 $\Gamma^k$：
\begin{equation}
  \Gamma^{k}_{\gamma\delta,\alpha\beta}(P_1,P_2)=\lim_{K\to0}\Gamma_{\gamma\delta,\alpha\beta}(K;P_1,P_2),\ \text{当}\ \omega/k\to0
\end{equation}
这个函数(乘以 $Z^2$)正是向前散射幅(即 $P_1,P_2\to P_1,P_2$ 的跃迁)，它相应于费米面上由准粒子产生的真实物理过程：即留在该表面的准粒子的碰撞引起无能量改变的动量变化，因此过渡到动量传递为零($k\to0$)的极限，应在能量的传递严格等于零($\omega=0$)的情况下进行.前面引入的函数 $\Gamma^{\omega}$，相应于动量传递严格等于零($k=0$)时小能量传递的非物理“散射”的极限情况.

在(17.15)式中，取 $\omega=0$，再过渡到 $\bm{k}=0$ 的极限，并以 $Z^2$ 乘等式的两边，则得：
\begin{equation}
\begin{aligned}
  Z^2\Gamma^{k}_{\gamma\delta,\alpha\beta}(P_1,P_2)&=Z^2\Gamma^{\omega}_{\gamma\delta,\alpha\beta}(P_1,P_2)-\\
  -{}&\frac{p_{\mathrm{F}}^2}{v_{\mathrm{F}}(2\pi)^3}\int Z^2\Gamma^{\omega}_{\gamma\zeta,\alpha\kappa}(P_1,Q_{\mathrm{F}})\cdot Z^2\Gamma^{k}_{\kappa\delta,\zeta\beta}(Q_{\mathrm{F}},P_2)\,\mathrm{d}o_l.
\end{aligned}
\end{equation}
因此，存在一个联系向前散射幅两种极限形式的普遍关系.

函数 $\Gamma$ 的反对称性质(15.8)，给出 $P_1\to P_2$ 时 $\Gamma^{k}$ 和 $\Gamma^{\omega}$ 行为的一个信息.在该等式中，取 $P_1=P_2$，$\alpha=\beta$，我们得到：
\begin{equation}
  \Gamma_{\gamma\delta,\alpha\alpha}(P_1+K,P_1-K;P_1,P_1)=0
\end{equation}
(这里不对 $\alpha$ 求和！)\footnote{当只考虑准粒子自旋之间的交换相互作用时，所有 $\Gamma_{\gamma\delta,\alpha\alpha}$ 中仅仅 $\Gamma_{\alpha\alpha,\alpha\alpha}$ 不等于零. 这个论断表明，散射时自旋矢量具有不变性. 此论断也可以直接根据(2.4)形的表达式来验证.}.在此等式中，应谨慎进行向 $\Gamma^{\omega}$ 或 $\Gamma^{k}$ 的过渡，因为在 $\Gamma^{\omega}$ 和 $\Gamma^{k}$ 中，首先取 $K=0$，而在(17.18)式中，则首先取 $P_1=P_2$.

同时令 $K$ 和 $P_1-P_2\equiv S=(s_0,\bm{s})$ 都是小量.这样，除图(17.2)以外，下图也将是危险的：
\begin{equation*}
  \setlength{\unitlength}{0.80pt}%
  \begin{picture}(310,112)(0,0)
    \put(58,94){\vector(1,-1){26}}%
    \put(232,94){\vector(-1,-1){26}}%
    \qbezier(95,51)(104.14,51.86)(105,61)%
    \qbezier(105,61)(104.14,70.14)(95,71)%
    \qbezier(95,71)(85.86,70.14)(85,61)%
    \qbezier(85,61)(85.86,51.86)(95,51)%
    \qbezier(195,51)(204.14,51.86)(205,61)%
    \qbezier(205,61)(204.14,70.14)(195,71)%
    \qbezier(195,71)(185.86,70.14)(185,61)%
    \qbezier(185,61)(185.86,51.86)(195,51)%
    \qbezier(95,71)(145,96)(195,71)%
    \put(145,104){\makebox(0,0){$Q{+}S{+}K$}}%
    \tarrL{195}{51}{94}%
    \put(150,36){\makebox(0,0){$Q$}}%
    \put(89,54){\vector(-1,-1){28}}%
    \put(201,54){\vector(1,-1){28}}%
    \put(44,104){\makebox(0,0){$P_1$}}%
    \put(246,104){\makebox(0,0){$P_2$}}%
    \put(42,14){\makebox(0,0){$P_2{-}K$}}%
    \put(250,14){\makebox(0,0){$P_1{+}K$}}%
  \end{picture}
\end{equation*}
因此，当 $K,S\to0$ 时，函数 $\Gamma_{\gamma\delta,\alpha\alpha}$ 将依赖于两个“特殊”宗量：
\begin{equation*}
  x=\frac{\omega}{k},\qquad y=\frac{s_0+\omega}{|\bm{s}+\bm{k}|}
\end{equation*}
当 $x=y$ 时，(17.18)式表示该函数变为零.我们来研究费米面上 $\Gamma$ 的值；这时 $\omega=s_0=0$，因而 $y=0$.所以只有当 $x=0$ 时，在这个极限上才有等式(17.18).换言之，在费米面上对于 $\Gamma^{k}$ 的等式是成立的：
\begin{equation}
  \Gamma^{k}_{\gamma\delta,\alpha\alpha}(P_1,P_1)=0
\end{equation}
(N. D. Mermin, 1967).

\section{顶角函数与准粒子相互作用函数的关系}

正如在定义单粒子格林函数的(7.9)式的矩阵元的结构中，包含粒子数为 $N\pm1$ 的中间态一样，在构成双粒子格林函数(矩阵元(15.1))中，也包含粒子数为 $N,N\pm1,N\pm2$ 的中间态\footnote{例如，粒子数为 $N$ 的状态在 $T$ 乘积的算符序列 $\hat{\Psi}_3\hat{\Psi}^{+}_1\hat{\Psi}_4\hat{\Psi}^{+}_2$ 中出现. 而粒子数为 $N+2$ 的状态，相应于 $\hat{\Psi}_3\hat{\Psi}_4\hat{\Psi}^{+}_1\hat{\Psi}^{+}_2$ 这样的序列.}.

由于存在粒子数为 $N\pm1$ 的中间态，双粒子格林函数的极点与 $G$ 函数的极点(即与准粒子能量)是一致的.但是，所对应的各因子，已在(15.7)式中以显形式被分解出来.因此，以这个顶角公式所定义的函数 $\Gamma$，只具有对应于粒子数为 $N$ 和 $N\pm2$ 的各态的极点.这些态的角动量与基态的角动量之差为 0 或 1，因此相应于这些极点的元激发具有整数自旋(0 或 1)，所以它们遵从玻色统计.换句话说，顶角函数的各极点决定费米液体能谱的玻色支.

由粒子数不变的中间态产生的极点，相应于元激发是零声量子.在图技术中，中间态相应于各图的不同断面，它们在图的某些外线之间将图分成两个部分.在这种情况下，粒子数不变的诸中间态相应于图(17.3)的各断面，它们分别截断连接相邻的 $\tilde{\Gamma}$ 单元的一对实线；在这些状态中，粒子数的不变性表现在穿过断面两侧的线数是相同的.通过这种截面所迁移的 4 维动量是：$(Q+K)-Q=K$；因此，粒子数不变的元激发则相应于变量为 $K$ 的顶角函数 $\Gamma(K;P_1,P_2)$ 的极点.

在前面［推导(17.10)式时］我们看到，在 4 维矢量 $Q$ 和 $Q+K$ 中的两个动量 $\bm q$ 和 $\bm q+\bm k$，其一应大于边界动量 $p_{\mathrm{F}}$，而另一个应小于 $p_{\mathrm{F}}$.从另一方面来看，当基态激发时，在费米球外的只能是“粒子”，而在球内的只能是“空穴”.就此意义可以说，费米液体中的零激发，可看作是粒子与空穴的束缚态\footnote{就问题的这种提法来说，在形式上与量子电动力学中(见第四卷 §125)电子和正电子束缚态能级的定义有许多共同之处. 特别是，方程(17.4—17.5)类似于第四卷的 Bethe--Salpeter 方程(125.10—125.11).}.

相应于 $N\pm2$ 个粒子中间态的元激发［它们相应于函数 $\Gamma(K;P_1,P_2)$ 对变量 $P_1+P_2$ 的极点］，可以看作两个粒子或两个空穴的束缚态.但是，这种状态的存在(第五章将要说明)，将使费米液体具有超流性，因而也需对图技术的全部数学工具作重要的改变.

因此，为了确定费米液体非超流能谱的玻色支，应当研究变量为 $K=(\omega,\bm{k})$ 的顶角函数 $\Gamma(K;P_1,P_2)$ 的极点.每取一个 $k$ 值，对应于极点便有一定的能量 $\omega=\omega(k)$，于是就确定了这些激发的色散律.对于弱激发态，$\omega$ 和 $k$ 都小，因此可以利用已得到的在 $K$ 为小值的区域内函数 $\Gamma(K;P_1,P_2)$ 的方程.

在函数 $\Gamma$ 的极点附近，方程(17.15)的左边及其右边的积分是非常大的量；而 $\Gamma^{\omega}(P_1,P_2)$ 项仍是有限的，因此可以略去.其次应注意，在方程(17.15)中对函数 $\Gamma$ 进行的运算没有涉及变量 $P_2$ 及下标 $\beta$ 和 $\delta$，即它们在方程中只起不重要的参量作用.最后，我们来研究在费米球面上的函数 $\Gamma$，即取 $P_1=(0,p_{\mathrm{F}}\bm{n})$，其中 $\bm n$ 为变单位矢量.基于所有上边这些论述，我们作出这样的结论：确定费米液体中的声激发，归结为求下列积分方程的本征值问题：
\begin{equation}
  \chi_{\gamma\alpha}(\bm{n})=\frac{Z^2p_{\mathrm{F}}^2}{(2\pi)^3}\int\Gamma^{\omega}_{\gamma\zeta,\alpha\kappa}(\bm{n},\bm{l})\chi_{\kappa\zeta}(\bm{l})\frac{\bm{l}\cdot\bm{k}\,\mathrm{d}o_l}{\omega-v_{\mathrm{F}}\bm{l}\cdot\bm{k}},
\end{equation}
式中 $\chi_{\gamma\alpha}(\bm n)$ 为辅助函数.

为改造这个方程，引入如下新的函数代替 $\chi$：
\begin{equation}
  \nu_{\gamma\alpha}(\bm{n})=\frac{\bm{n}\cdot\bm{k}}{\omega-v_{\mathrm{F}}\bm{n}\cdot\bm{k}}\chi_{\gamma\alpha}(\bm{n}),
\end{equation}
这时，方程(18.1)取如下形式：
\begin{equation}
  \bigl(\omega-v_{\mathrm{F}}\bm{n}\cdot\bm{k}\bigr)\nu_{\gamma\alpha}(\bm{n})=\bm{k}\cdot\bm{n}\,\frac{p_{\mathrm{F}}^2Z^2}{(2\pi)^3}\int\Gamma^{\omega}_{\gamma\zeta,\alpha\kappa}(\bm{n},\bm{n}')\nu_{\kappa\zeta}(\bm{n}')\,\mathrm{d}o'
\end{equation}
(记号 $\bm l$ 换成 $\bm n'$).

这个方程的形状，同费米液体振动的动理学方程(4.10)完全一致.将两个方程作对比，可导出准粒子相互作用函数与函数 $\Gamma^{\omega}$ 之间的下列对应关系：
\begin{equation}
  f_{\gamma\delta,\alpha\beta}(p_{\mathrm{F}}\bm{n},p_{\mathrm{F}}\bm{n}')=Z^2\Gamma^{\omega}_{\gamma\delta,\alpha\beta}(\bm{n},\bm{n}').
\end{equation}
因而弄清楚了函数 $f$ 和准粒子散射性质之间的关系\footnote{上述的一般推导属于 Л. Д. 朗道的工作(1958). 稍早一些，А. Б. Мигдал 和 В. М. Галицкий 对具体的(17.3)型图进行求和，从而得出弱非理想费米气体的动理学方程. 应当指出，在气体情况下，$G$ 函数中(在零级近似)只含有极点项，因而不存在消除非极点项的问题.}.

等式(18.4)将 $f$ 与非物理过程散射幅联系起来了.现在利用公式(17.17)得出 $f$ 同费米面上准粒子的“物理”向前散射幅之间的显示关系，我们把该散射幅表示成：
\begin{equation}
  A_{\gamma\delta,\alpha\beta}(\bm{n}_1,\bm{n}_2)=Z^2\Gamma^{k}_{\gamma\delta,\alpha\beta}(\bm{n}_1,\bm{n}_2).
\end{equation}
在费米面上，关系式(17.17)取如下形式：
\begin{equation}
\begin{aligned}
  &A_{\gamma\delta,\alpha\beta}(\bm{n}_1,\bm{n}_2)=\\
  ={}&f_{\gamma\delta,\alpha\beta}(\bm{n}_1,\bm{n}_2)-\frac{p_{\mathrm{F}}^2}{(2\pi)^2v_{\mathrm{F}}}\int f_{\gamma\zeta,\alpha\chi}(\bm{n}_1,\bm{n}')A_{\chi\delta,\zeta\beta}(\bm{n}',\bm{n}_2)\frac{\mathrm{d}o'}{4\pi}.
\end{aligned}
\end{equation}
函数 $A$ 和 $f$ 的自旋关系可以用泡利矩阵 $\bm{\sigma}$ 来表达.在一般情况下，这两个函数可以包含四个矢量 $\bm{n}_1,\bm{n}_2,\bm{\sigma}_1,\bm{\sigma}_2$ 的任何标量组合.但如果粒子间的作用是交换相互作用，则可容许的标积只是 $\bm{n}_1\cdot\bm{n}_2$ 和 $\bm{\sigma}_1\cdot\bm{\sigma}_2$.这样，函数 $A$ 和 $f$［如在(2.4)式中对待 $f$ 的作法一样］可以表为：
\begin{equation}
\begin{aligned}
  \frac{p_{\mathrm{F}}^2}{\pi^2v_{\mathrm{F}}}f_{\gamma\delta,\alpha\beta}(\bm{n}_1,\bm{n}_2)&=F(\vartheta)\delta_{\alpha\gamma}\delta_{\beta\delta}+G(\vartheta)\bm{\sigma}_{\gamma\alpha}\cdot\bm{\sigma}_{\delta\beta},\\
  \frac{p_{\mathrm{F}}^2}{\pi^2v_{\mathrm{F}}}A_{\gamma\delta,\alpha\beta}(\bm{n}_1,\bm{n}_2)&=B(\vartheta)\delta_{\alpha\gamma}\delta_{\beta\delta}+G(\vartheta)\bm{\sigma}_{\gamma\alpha}\cdot\bm{\sigma}_{\delta\beta},
\end{aligned}
\end{equation}
其中系数 $F,G,B,C$ 只是 $\bm n_1$ 和 $\bm n_2$ 之间夹角 $\vartheta$ 的函数.将这些函数按勒让德多项式展开：
\begin{equation}
  B(\vartheta)=\sum_{l=0}^{\infty}\,(2l+1)B_lP_l(\cos\vartheta),\cdots,
\end{equation}
将(18.7—18.8)式代入(18.6)式，并算出积分(这时利用勒让德多项式加法定理)，可得：
\begin{equation}
  B_l=F_l(1-B_l),\quad C_l=G_l(1-C_l).
\end{equation}
这些公式建立了 $f$ 和 $A$ 的展开系数间的简单代数关系.

由稳定性条件(2.19—2.20)可导出类似的关于系数 $B_l,C_l$ 的不等式：
\begin{equation}
  B_l<1,\quad C_l<1.
\end{equation}
此外，这些系数还满足由公式(17.19)得出的关系式：
\begin{equation}
\begin{gathered}
  B(0)+C(0)=0\\
  \text{或}\qquad\sum_{l=0}^{\infty}\,(2l+1)(B_l+C_l)=0.
\end{gathered}
\end{equation}
等式(18.9)，(18.11)连同条件(18.10)足以证明一个有趣的论断：在一切稳定的费米液体中，至少有一个(寻常的或有自旋的)轴对称的零声支\footnote{见 N. D. Mermin, \emph{Phys. Rev.} \textbf{159}, 161 (1967).}.

\section{格林函数微商恒等式}

在有关格林函数的数学工具中，格林函数的微商和准粒子散射幅之间的某些恒等关系起着重要作用.这些关系式的推导都是同样的，即算出在某一虚设的“外场”影响下格林函数的改变，而该外场对系统作用的结果预先已是已知的.

因此，首先要算出在任意“外场”影响下格林函数的改变 $\delta G$.这种场在哈密顿量中的对应项为：
\begin{equation}
  \delta\hat{V}^{(1)}=\int\hat{\Psi}^{*}_{\alpha}(t,\bm{r})\,\delta\hat{U}\,\hat{\Psi}_{\alpha}(t,\bm{r})\,\mathrm{d}^3x,
\end{equation}
其中 $\delta\hat{U}$ 为某一算符，它作用在 $\bm r$(也可以和时间 $t$ 有关)的函数上.

当有外场存在时，格林函数已与两个 4 维动量 $P_1$ 和 $P_2$ 有关.在图技术中，这种外场可用新的图元素——外虚线来表示：
\begin{equation*}
  \setlength{\unitlength}{0.90pt}%
  \begin{picture}(240,70)(0,0)
    \tarrL{200}{35}{120}%
    \dvh{140}{39}{4}%
    \put(64,48){\makebox(0,0){$P_2$}}%
    \put(214,48){\makebox(0,0){$P_1$}}%
  \end{picture}
\end{equation*}
并且这种线相当于因子
\begin{equation}
  -\mathrm{i}\delta U(P_2,P_1)=-\mathrm{i}\int\mathrm{e}^{\mathrm{i}P_2X}\delta\hat{U}\,\mathrm{e}^{-\mathrm{i}P_1X}\,\mathrm{d}^4X.
\end{equation}
外场对于精确格林函数的第一级修正可表为两个骨架图之和：
\begin{equation}
  \setlength{\unitlength}{0.75pt}%
  \begin{picture}(500,130)(0,0)
    \put(66,64){\makebox(0,0){$\mathrm{i}\delta G(P_2,P_1)=$}}%
    \tarrL{250}{64}{100}%
    \dvh{200}{68}{4}%
    \put(134,48){\makebox(0,0){$P_2$}}%
    \put(264,48){\makebox(0,0){$P_1$}}%
    \put(286,64){\makebox(0,0){$+$}}%
    \put(356,116){\vector(-1,-1){28}}%
    \qbezier(324,71)(335.882,72.118)(337,84)%
    \qbezier(337,84)(335.882,95.882)(324,97)%
    \qbezier(324,97)(312.118,95.882)(311,84)%
    \qbezier(311,84)(312.118,72.118)(324,71)%
    \put(312,72){\vector(-1,-1){32}}%
    \qbezier(334,90)(372,112)(410,90)%
    \qbezier(334,78)(372,58)(410,78)%
    \put(368,109){\vector(1,0){12}}%
    \put(376,60){\vector(-1,0){12}}%
    \put(372,122){\makebox(0,0){$Q_2$}}%
    \put(372,46){\makebox(0,0){$Q_1$}}%
    \multiput(412,80)(7,0){4}{\line(1,0){5}}%
  \end{picture}
\end{equation}
其中所有实线——粗线为精确 $G$ 函数，小圆为精确顶角函数($\mathrm{i}\Gamma$).将该等式以解析形式写出来，则有
\begin{equation}
\begin{aligned}
  \delta G_{\beta\alpha}(P_2,P_1)&=G_{\beta\gamma}(P_2)\delta U(P_2,P_1)G_{\gamma\alpha}(P_1)-\\
  -\mathrm{i}G_{\beta\gamma}&(P_2)G_{\varepsilon\alpha}(P_1)\int\Gamma_{\gamma\delta,\varepsilon\zeta}(P_2,Q_1;P_1,Q_2)\times\\
  &\times\delta U(Q_2,Q_1)G_{\zeta\kappa}(Q_2)G_{\kappa\delta}(Q_1)\frac{\mathrm{d}^4Q_1}{(2\pi)^4},
\end{aligned}
\end{equation}
并且 $Q_2+P_1=P_2+Q_1$.

我们感兴趣的头两个恒等式与系统中粒子数守恒有关.这一性质表现在系统哈密顿量中 $\psi$ 算符是成对出现的，即对每个宗量 $X$，算符 $\hat{\Psi}^{+}(X)$ 和 $\hat{\Psi}(X)$ 各出现一次.

现在对 $\psi$ 算符进行规范变换：
\begin{equation}
  \hat{\Psi}_{\alpha}(X)=\hat{\Psi}'_{\alpha}(X)\mathrm{e}^{-\mathrm{i}\chi(X)},\qquad
  \hat{\Psi}^{+}_{\alpha}=\hat{\Psi}^{\prime+}_{\alpha}\,\mathrm{e}^{\mathrm{i}\chi(X)},
\end{equation}
式中 $\chi(X)$ 为实函数\footnote{这类似于量子电动力学中的规范变换(见第三卷)(111.8—111.9).}.根据前面指出的哈密顿量的性质，如果 $\hat{\Psi}$ 遵从“薛定谔方程”(7.8)，$\hat{\Psi}'$ 则遵从作如下代换的薛定谔方程：
\begin{equation*}
  \Delta\to(\nabla-\mathrm{i}\,\nabla\chi)^2,\qquad
  \frac{\partial}{\partial t}\longrightarrow\frac{\partial}{\partial t}-\mathrm{i}\,\frac{\partial\chi}{\partial t}.
\end{equation*}
当 $\chi=\delta\chi$ 为无穷小量时，方程的这个改变相当于给哈密顿量附加一“外场”：
\begin{equation*}
  \delta\hat{U}=-\frac{\partial\delta\chi}{\partial t}+\frac{\mathrm{i}}{2m}\bigl(\Delta\delta\chi+2(\nabla\delta\chi)\nabla\bigr).
\end{equation*}
特别是，如果
\begin{equation*}
  \delta\chi(X)=\mathrm{Re}(\chi_0\mathrm{e}^{-\mathrm{i}KX}),\qquad K=(\omega,\bm{k})
\end{equation*}
(由于以后的运算具有线性，因此记号 Re 可以略去)，所以
\begin{equation}
  \delta U(P_2,P_1)=\mathrm{i}(2\pi)^4\chi_0\delta^{(4)}(P_2-P_1-K)\biggl\{\omega-\frac{1}{2m}\bm{k}(\bm{p}_1+\bm{p}_2)\biggr\}.
\end{equation}
另一方面，由 $\psi$ 算符
\begin{equation*}
  \hat{\Psi}'_{\alpha}=\hat{\Psi}_{\alpha}(1+\mathrm{i}\delta\chi),\qquad
  \hat{\Psi}^{\prime+}_{\alpha}=\hat{\Psi}^{+}_{\alpha}(1-\mathrm{i}\delta\chi)
\end{equation*}
构成的格林函数与由算符 $\hat{\Psi},\hat{\Psi}^{+}$ 构成的函数之差为：
\begin{equation*}
  \delta G_{\alpha\beta}(X_1,X_2)=\mathrm{i}G_{\alpha\beta}(X_1,X_2)\bigl[\delta\chi(X_1)-\delta\chi(X_2)\bigr],
\end{equation*}
或以傅里叶分量表示：
\begin{equation}
\begin{aligned}
  \delta G_{\alpha\beta}(P_2,&P_1)=\int\delta G_{\alpha\beta}(X_1,X_2)\mathrm{e}^{\mathrm{i}(P_2X_1-P_1X_2)}\mathrm{d}^4X_1\mathrm{d}^4X_2=\\
  ={}&\mathrm{i}\bigl[G_{\alpha\beta}(P_1)-G_{\alpha\beta}(P_2)\bigr]\delta\chi(P_2-P_1),
\end{aligned}
\end{equation}
式中
\begin{equation*}
  \delta\chi(P)=\int\delta\chi(X)\mathrm{e}^{\mathrm{i}PX}\mathrm{d}^4X=(2\pi)^4\chi_0\delta^{(4)}(P-K).
\end{equation*}
因此，同一个变化 $\delta G_{\alpha\beta}$ 可表为两种形式：(19.7)式和(19.4)式［需要将(19.6)式的 $\delta U$ 代入后者).使这两个表达式相等，(经过 $G_{\alpha\beta}=G\delta_{\alpha\beta}$ 代换和重新表示某些变量之后)我们得到：
\begin{equation*}
\begin{aligned}
  \delta_{\alpha\beta}\bigl[G(P+K)-G(P)\bigr]=G(P+K)G(P)\biggl\{\biggl[-\omega+\frac{\bm{k}\cdot(2\bm{p}+\bm{k})}{2m}\biggr]\delta_{\alpha\beta}+\\
  +\mathrm{i}\int\Gamma_{\beta\delta,\alpha\delta}(K;P,Q)G(Q)G(Q-K)\biggl[\omega-\frac{\bm{k}\cdot(2\bm{q}-\bm{k})}{2m}\biggr]\frac{\mathrm{d}^4Q}{(2\pi)^4}\biggr\}.
\end{aligned}
\end{equation*}
在这个等式中取 $\omega,\bm{k}\to0$ 时的极限，便得到待求的恒等式；这时
\begin{equation}
  G(P+K)-G(P)\longrightarrow\omega\frac{\partial G}{\partial p_0}+\bm{k}\frac{\partial G}{\partial\bm{p}}
\end{equation}
［式中 $P=(p_0,\bm{p})$］.在 $k/\omega\to0$ 的条件下取此极限，我们得到第一个恒等式：
\begin{equation}
  \delta_{\alpha\beta}\frac{\partial G(P)}{\partial p_0}
  =-\bigl\{G^2(P)\bigr\}_{\omega}\biggl[\delta_{\alpha\beta}-\mathrm{i}\int\Gamma^{\omega}_{\beta\delta,\alpha\delta}(P,Q)\bigl\{G^2(Q)\bigr\}_{\omega}\frac{\mathrm{d}^4Q}{(2\pi)^4}\biggr],
\end{equation}
这里引入了记号
\begin{equation}
  \bigl\{G^2(P)\bigr\}_{\omega}=\lim_{\omega,k\to0}G(P)G(P+K),\ \text{当}\ k/\omega\to0\ \text{时}.
\end{equation}
用类似方法，在 $\omega/k\to0$ 的条件下取极限，我们又得到一个恒等式：
\begin{equation}
  \delta_{\alpha\beta}\frac{\partial G}{\partial\bm{p}}
  =\bigl\{G^2(P)\bigr\}_{k}\biggl[\frac{\bm{p}}{m}\delta_{\alpha\beta}-\mathrm{i}\int\Gamma^{k}_{\beta\delta,\alpha\delta}(P,Q)\frac{\bm{q}}{m}\bigl\{G^2(Q)\bigr\}_{k}\frac{\mathrm{d}^4Q}{(2\pi)^4}\biggr],
\end{equation}
其中引入了类似的记号 $\{G^2(P)\}_k$.

下面来研究给系统加上如下恒定场时格林函数的改变：
\begin{equation}
  \delta\hat{U}=\delta U(\bm{r})=U_0\mathrm{e}^{\mathrm{i}\bm{k}\cdot\bm{r}}.
\end{equation}
当 $\bm{k}\to0$ 时，这个场在空间缓慢地改变，因此可以把它视为对系统的宏观影响.根据外场中的热力学平衡条件，应有 $\mu+\delta U=$ 常量(见第五卷 §25)；这就是说，当 $\bm{k}\to0$ 时化学势 $\mu$ 改变一个小量 $-U_0$.格林函数相应的改变为：
\begin{equation*}
  \delta G_{\alpha\beta}(X_1,X_2)=-U_0\delta_{\alpha\beta}\frac{\partial G(X_1-X_2)}{\partial\mu}.
\end{equation*}
它的傅里叶分量［按(19.7)式的定义］：
\begin{equation*}
  \delta G_{\alpha\beta}(P_2,P_1)=-(2\pi)^4\delta^{(4)}(P_2-P_1)U_0\frac{\partial G(P_1)}{\partial\mu}.
\end{equation*}
另一方面，格林函数的这个改变也可以按公式(19.4)计算，这一次在公式中取
\begin{equation*}
  \delta U(P_2,P_1)=(2\pi)^4U_0\,\delta^{(4)}(P_2-P_1-K),\ (\ K=0,\bm{k}\ ).
\end{equation*}
在此情况下(恒定场，$\omega\equiv0$)，向 $\bm{k}\to0$ 的极限过渡相应于 $\omega/k\to0$ 的情形.结果得到恒等式：
\begin{equation}
  \delta_{\alpha\beta}\frac{\partial G(P)}{\partial\mu}
  =-\bigl\{G^2(P)\bigr\}_{k}\biggl[\delta_{\alpha\beta}-\mathrm{i}\int\Gamma^{k}_{\beta\delta,\alpha\delta}(P,Q)\bigl\{G^2(Q)\bigr\}_{k}\frac{\mathrm{d}^4Q}{(2\pi)^4}\biggr].
\end{equation}
最后，这一恒等式的出现，可作为系统伽利略不变性的结果.为导出该恒等式，我们来研究动坐标系中的液体，此坐标系以微小速度 $\delta\bm{w}(t)=\bm{w}_0\mathrm{e}^{-\mathrm{i}\omega t}$ 随时间作缓慢运动.变换到这个坐标系，等价于给系统加上外场，外场的算符为\footnote{变换到动坐标系时，需要在自由粒子经典拉格朗日函数 $L=mv^2/2$ 中作代换：$\bm{v}\to\bm{v}+\delta\bm{w}$，于是出现一个小增量（当 $\delta\bm{w}$ 很小时）$\delta L=m\bm{v}\cdot\delta\bm{w}$(当 $\delta\bm{w}$ 很小时). 相应地［见第一卷(40.7)］，哈密顿函数增量为 $\delta H=-\bm{p}\cdot\delta\bm{w}$，在量子力学中该增量相应于算符(19.14).}：
\begin{equation}
  \delta\hat{U}=-\delta\bm{w}\cdot\bm{p}=\frac{\mathrm{i}}{m}\delta\bm{w}\cdot\nabla.
\end{equation}
或在动量表象中，
\begin{equation*}
  \delta U(P_2,P_1)=-\bm{p}_i\cdot\bm{w}_0(2\pi)^4\delta^{(4)}(P_2-P_1-K),\quad K=(\omega,0),
\end{equation*}
此表达式必须代入(19.4)式，然后取极限 $\omega\to0$.

另一方面，当 $\omega\to0$ 时，这里所指的是从一个惯性参考系向以恒定速度 $\delta\bm w$ 运动的另一惯性参考系作伽利略变换.如果液体有能量为 $\varepsilon(\bm{p})$ 的元激发，那么在以速度 $\delta\bm w$ 相对于液体运动的参考系中，这个元激发的能量将为 $\varepsilon-\bm{p}\cdot\delta\bm{w}$\footnote{对照下面 §23 中更详细的讨论.}.因此在新的参考系中，频率 $p_0$ 应以 $p_0+\bm{p}\cdot\delta\bm{w}$ 的组合形式出现在函数 $G(F)$ 中(于是，使函数的极点移动 $-\bm{p}\cdot\delta\bm{w}$).这样一来，
\begin{equation*}
  \delta G=\bm{p}\cdot\delta\bm{w}\,\frac{\partial G}{\partial p_0},
\end{equation*}
于是我们求得恒等式：
\begin{equation}
  \delta_{\alpha\beta}\bm{p}\,\frac{\partial G(P)}{\partial p_0}
  =-\bigl\{G^2(P)\bigr\}_{\omega}\biggl\{\delta_{\alpha\beta}\bm{p}-\mathrm{i}\int\Gamma^{\omega}_{\beta\delta,\alpha\delta}(P,Q)\bm{q}\bigl\{G^2(Q)\bigr\}_{\omega}\frac{\mathrm{d}^4Q}{(2\pi)^4}
\end{equation}
下面我们要运用得到的恒等式，特别是，运用于自由变量 $P=(p_0,\bm{p})$ 在费米面上的值为 $P_{\mathrm{F}}=(0,p_{\mathrm{F}})$ 的情形.将因子 $G^2(P)$ 从恒等式的右边移至左边，同时把 $G(P)$ 的微商换成 $G^{-1}(P)$ 的微商；这时在 $G(P)G(P+K)$ 中以什么方式趋于 $K\to0$ 的极限便是无关紧要的了.

另一方面，在费米面附近格林函数决定于本身的极点项，因此
\begin{equation*}
  G^{-1}(P)=\frac{1}{Z}\bigl[p_0-v_{\mathrm{F}}(p-p_{\mathrm{F}})\bigr],
\end{equation*}
从而，在该表面上有：
\begin{equation*}
  \frac{\partial G^{-1}}{\partial p_0}=\frac{1}{Z},\qquad
  \frac{\partial G^{-1}}{\partial\mu}=\frac{v_{\mathrm{F}}}{Z}\,\frac{\mathrm{d}p_{\mathrm{F}}}{\mathrm{d}\mu}.
\end{equation*}
结果，例如恒等式(19.9)和(19.13)在费米面上分别取如下形式：
\begin{equation}
  \mathrm{i}\int\Gamma^{\omega}_{\beta\delta,\alpha\delta}(P_{\mathrm{F}},Q)\bigl\{G^2(Q)\bigr\}_{\omega}\frac{\mathrm{d}^4Q}{(2\pi)^4}=\Bigl(1-\frac{1}{Z}\Bigr)\delta_{\alpha\beta},
\end{equation}
\begin{equation}
  \mathrm{i}\int\Gamma^{k}_{\beta\delta,\alpha\delta}(P_{\mathrm{F}},Q)\bigl\{G^2(Q)\bigr\}_{k}\frac{\mathrm{d}^4Q}{(2\pi)^4}=\Bigl(1-\frac{v_{\mathrm{F}}}{Z}\frac{\mathrm{d}p_{\mathrm{F}}}{\mathrm{d}\mu}\Bigr)\delta_{\alpha\beta}.
\end{equation}

\section{边界动量与密度关系的推导}

前几节得到的诸关系式，对朗道费米液体理论的基本命题能够给出彻底的证明：即断定边界动量 $p_{\mathrm{F}}$ 与液体密度 $N/V$ 之间的关系，可由适用于理想气体的同样公式(1.1)给出.

这里证明的主旨在于，当化学势 $\mu$ 的变化为无穷小时，可独立计算 $N$ 和 $p_{\mathrm{F}}$ 的改变，然后再将它们作对比.

根据(7.24)式，作为(在给定的体积 $V$ 中)化学势函数的总粒子数，可由下列积分给出：
\begin{equation}
  N=-2\mathrm{i}V\lim_{t\to-0}\int G(P)\mathrm{e}^{-\mathrm{i}p_0t}\frac{\mathrm{d}^4p}{(2\pi)^4},\qquad P=(p_0,\bm{p}).
\end{equation}
因而微商
\begin{equation}
  \frac{1}{V}\frac{\mathrm{d}N}{\mathrm{d}\mu}=-2\mathrm{i}\int\frac{\partial G(P)}{\partial\mu}\frac{\mathrm{d}^4P}{(2\pi)^4}.
\end{equation}
由于这个积分当 $p_0$ 大时($|p_0|\to\infty$ 时，$\partial G/\partial\mu\propto1/p_0^2$)具有收敛性，因此在被积式中已不需要写出因子 $\mathrm{e}^{-\mathrm{i}p_0t}$.将恒等式(19.13)(对 $\alpha=\beta$ 的值求和)中的 $\partial G/\partial\mu$ 代入积分后，我们求得：
\begin{equation*}
  \frac{1}{V}\frac{\mathrm{d}N}{\mathrm{d}\mu}=2\mathrm{i}\int\bigl\{G^2(P)\bigr\}_{k}\frac{\mathrm{d}^4P}{(2\pi)^4}
  +\int\bigl\{G^2(P)\bigr\}_{k}\Gamma^{k}(P,Q)\bigl\{G^2(Q)\bigr\}_{k}\frac{\mathrm{d}^4P\mathrm{d}^4Q}{(2\pi)^8}.
\end{equation*}
为简便起见，式中 $\Gamma=\Gamma_{\alpha\gamma,\alpha\gamma}$.以下计算的目的，是将等式的右边部分只用沿费米面的积分表达出来.

首先，用(17.17)的表达式(其中将记号 $Q_{\mathrm{F}}$ 换成 $S_{\mathrm{F}}$)代替上面的第二个积分中的 $\Gamma^{k}$；
\begin{equation}
\begin{aligned}
  \frac{1}{V}\frac{\mathrm{d}N}{\mathrm{d}\mu}={}&2\mathrm{i}\int\bigl\{G^2(P)\bigr\}_{k}\frac{\mathrm{d}^4P}{(2\pi)^4}
  +\int\bigl\{G^2(P)\bigr\}_{k}\Gamma^{\omega}(P,Q)\bigl\{G^2(Q)\bigr\}_{k}\frac{\mathrm{d}^4P\mathrm{d}^4Q}{(2\pi)^8}-\\
  -\frac{p_{\mathrm{F}}^2Z^2}{v_{\mathrm{F}}(2\pi)^3}&\int\bigl\{G^2(P)\bigr\}_{k}\Gamma^{\omega}_{\alpha\zeta,\alpha\kappa}(P,S_{\mathrm{F}})\Gamma^{k}_{\kappa\gamma,\zeta\gamma}(S_{\mathrm{F}},Q)\bigl\{G^2(Q)\bigr\}_{k}\frac{\mathrm{d}^4P\mathrm{d}^4Q\mathrm{d}o_S}{(2\pi)^8}.
\end{aligned}
\end{equation}
我们首先变换最后一项.这一项的被积式中，只最后两个因子与 $Q$ 有关；它们对 $\mathrm{d}^4Q$ 的积分决定于公式(19.17)(在费米面上，$S=S_{\mathrm{F}}$)，因此这一项取如下形式：
\begin{equation*}
  \mathrm{i}\,\frac{p_{\mathrm{F}}^2Z^2}{v_{\mathrm{F}}(2\pi)^3}\int\bigl\{G^2(P)\bigr\}_{k}\Gamma^{\omega}(P,S_{\mathrm{F}})
  \frac{\mathrm{d}^4P\mathrm{d}o_S}{(2\pi)^4}\Bigl(1-\frac{v_{\mathrm{F}}}{Z}\frac{\mathrm{d}p_{\mathrm{F}}}{\mathrm{d}\mu}\Bigr).
\end{equation*}
其次我们记起，对 $\mathrm{d}^4P$ 取积分时，$G(P)G(P+K)$ 的极限值应按(17.10)式的意义来理解；所以 $\{G^2(P)\}_{\omega}=\varphi(P)$，而
\begin{equation}
  \bigl\{G^2(P)\bigr\}_{k}=\bigl\{G^2(P)\bigr\}_{\omega}-\frac{2\pi\mathrm{i}Z^2}{v_{\mathrm{F}}}\delta(p_0)\delta(p-p_{\mathrm{F}}).
\end{equation}
作此代换之后，则得：
\begin{equation*}
  \mathrm{i}\,\frac{p_{\mathrm{F}}^2Z^2}{v_{\mathrm{F}}(2\pi)^3}\Bigl(1-\frac{v_{\mathrm{F}}}{Z}\frac{\mathrm{d}p_{\mathrm{F}}}{\mathrm{d}\mu}\Bigr)
  \biggl\{\int\bigl\{G^2(P)\bigr\}_{\omega}\Gamma^{\omega}(P,S_{\mathrm{F}})\frac{\mathrm{d}^4P\mathrm{d}o_S}{(2\pi)^4}-8\pi\mathrm{i}\,\overline{F}\biggr\},
\end{equation*}
其中，根据(18.4)式引入了准粒子相互作用函数，并利用了(2.6—2.7)式中由函数 $F(\vartheta)$ 表示的表达式 $f_{\alpha\xi,\alpha\xi}$；$F$ 上面的横线表示对 $\mathrm{d}o/4\pi$ 求积分.剩下的对 $\mathrm{d}^4P$ 的积分，可由公式(19.16)给出，然后对 $\mathrm{d}o_S$ 求积分又得出因子 $4\pi$.因而(20.3)式中的第三项等于：
\begin{equation}
  -\frac{p_{\mathrm{F}}^2Z^2}{v_{\mathrm{F}}\pi^2}\biggl(\frac{v_{\mathrm{F}}}{Z}\frac{\mathrm{d}p_{\mathrm{F}}}{\mathrm{d}\mu}-1\biggr)\biggl\{1-\frac{1}{Z}+\overline{F}\biggr\}.
\end{equation}
用类似的方法来变换(20.3)式中的第二项：先根据(20.4)式，用 $\{G^2(P)\}_{\omega}$ 和 $\{G^2(Q)\}_{\omega}$ 表达 $\{G^2(P)\}_{k}$ 和 $\{G^2(Q)\}_{k}$，此后再运用恒等式(19.9)和(19.16).因而这一项等于
\begin{equation}
  -2\mathrm{i}\int\frac{\partial G}{\partial p_0}\frac{\mathrm{d}^4P}{(2\pi)^4}
  -2\mathrm{i}\int\bigl\{G^2(P)\bigr\}_{\omega}\frac{\mathrm{d}^4P}{(2\pi)^4}
  +\frac{p_{\mathrm{F}}^2Z^2}{v_{\mathrm{F}}\pi^2}\biggl\{2\Bigl(\frac{1}{Z}-1\Bigr)-\overline{F}\biggr\}.
\end{equation}
由于 $p_0\to\pm\infty$ 时 $G\to0$，所以对 $\mathrm{d}p_0$ 求积分时，第一个积分变为零.

最后，(20.3)式的第一项因代入(20.4)式后，而得出
\begin{equation}
  2\mathrm{i}\int\bigl\{G^2(P)\bigr\}_{\omega}\frac{\mathrm{d}^4P}{(2\pi)^4}+\frac{p_{\mathrm{F}}^2}{v_{\mathrm{F}}}\frac{Z^2}{\pi^2}.
\end{equation}
现在把所有的贡献(20.5—20.7)都加起来，我们求得：
\begin{equation}
  \frac{1}{V}\frac{\mathrm{d}N}{\mathrm{d}\mu}=\frac{p_{\mathrm{F}}^2}{\pi^2}\frac{\mathrm{d}p_{\mathrm{F}}}{\mathrm{d}\mu}
  +\frac{p_{\mathrm{F}}^2Z}{\pi^2v_{\mathrm{F}}}\biggl\{1-\frac{\mathrm{d}p_{\mathrm{F}}}{\mathrm{d}\mu}v_{\mathrm{F}}\bigl(1+\overline{F}\bigr)\biggr\}.
\end{equation}
另一方面，在(2.14)式中取
\begin{equation*}
  \delta n'=\frac{\partial n'}{\partial p_{\mathrm{F}}}\,\delta p_{\mathrm{F}}=\delta(p-p_{\mathrm{F}})\,\delta p_{\mathrm{F}},
\end{equation*}
得出：
\begin{equation}
  \frac{\mathrm{d}\mu}{\mathrm{d}p_{\mathrm{F}}}=v_{\mathrm{F}}\bigl(1+\overline{F}\bigr),
\end{equation}
我们强调，在推导(2.14)式时还没利用 $p_{\mathrm{F}}$ 与 $N/V$ 的具体依赖关系，因此我们有权在这里利用此关系式［指(2.14)—译注］来寻求上述依赖关系［当然，等式(20.9)也可以借助于推导(20.8)式时用过的那些关于顶角函数的诸关系式得出\footnote{有效质量公式(2.11)可借助关系式(17.17)以及等式(19.11)和(19.15)推导出来.}］.

考虑上面的等式我们看出，(20.8)式中的花括号变为零，所以
\begin{equation}
  \frac{\mathrm{d}}{\mathrm{d}\mu}\frac{N}{V}=\frac{p_{\mathrm{F}}^2}{\pi^2}\frac{\mathrm{d}p_{\mathrm{F}}}{\mathrm{d}\mu}
  =\frac{\mathrm{d}}{\mathrm{d}\mu}\biggl[\frac{8\pi p_{\mathrm{F}}^3}{3(2\pi)^3}\biggr].
\end{equation}
当 $N/V\to0$ 时，我们研究的对象则是气体，因此在这个极限下，$p_{\mathrm{F}}$ 与 $N/V$ 的依赖关系不管怎样都应与气体的一致.用这个条件可确定积分(20.10)式时出现的常数，于是最终我们得出了所求的关系式(1.1)：
\begin{equation*}
  \frac{N}{V}=\frac{8\pi p_{\mathrm{F}}^3}{3(2\pi)^3}.
\end{equation*}

\section{近理想费米气体格林函数}

为了举例说明图技术的应用方法，在 §6 中用通常的微扰论已经研究过［В. М. Галицкий, 1958］的模型框架内，本节将用图技术计算近理想费米气体的格林函数.注意这里所说的气体，是指粒子间具有斥力作用的气体，只要最终计算结果仅含散射幅，则 §6 中描写的方法便允许将微扰论运用于这种相互作用.

§14 指出过，求格林函数归结为计算自能函数 $\Sigma_{\alpha\beta}(P)$.在微扰论的一级和二级近似下，自能函数由(14.9)和(14.10)图的集合给出.在此，把它们画成如下形状：
\begin{equation}
  \setlength{\unitlength}{0.70pt}%
  \begin{picture}(640,195)(0,0)
    \put(14,160){\makebox(0,0){$P$}}%
    \arrL{56}{148}{36}%
    \qbezier(84,128)(102.28,129.72)(104,148)%
    \qbezier(104,148)(102.28,166.28)(84,168)%
    \qbezier(84,168)(65.72,166.28)(64,148)%
    \qbezier(64,148)(65.72,129.72)(84,128)%
    \put(84,148){\makebox(0,0){$-\mathrm{i}\Sigma$}}%
    \arrL{146}{148}{38}%
    \put(158,160){\makebox(0,0){$P$}}%
    \put(178,148){\makebox(0,0){$=$}}%
    \put(240,168){\thicklines\line(1,0){70}}%
    \qbezier(240,168)(216,168)(216,148)%
    \qbezier(216,148)(216,128)(240,128)%
    \put(240,128){\thicklines\line(1,0){70}}%
    \arrL{366}{168}{54}%
    \put(378,180){\makebox(0,0){$P$}}%
    \arrL{308}{156}{52}%
    \put(246,144){\makebox(0,0){$P$}}%
    \dvh{280}{130}{4}%
    \put(292,142){\makebox(0,0){$Q$}}%
    \arrL{372}{128}{58}%
    \put(384,128){\makebox(0,0){$P$}}%
    \put(275,184){\makebox(0,0){$P{-}Q$}}%
    \put(280,108){\makebox(0,0){(a)}}%
    \put(386,148){\makebox(0,0){$+$}}%
    \put(430,170){\thicklines\line(1,0){60}}%
    \qbezier(430,170)(418,170)(418,158)%
    \qbezier(418,158)(418,146)(430,146)%
    \qbezier(490,170)(502,170)(502,158)%
    \qbezier(502,158)(502,146)(490,146)%
    \tarrL{490}{146}{50}%
    \dvh{460}{104}{6}%
    \arrL{516}{102}{52}%
    \arrL{566}{102}{48}%
    \put(408,92){\makebox(0,0){$P$}}%
    \put(576,92){\makebox(0,0){$P$}}%
    \put(460,82){\makebox(0,0){(b)}}%
    \put(244,60){\makebox(0,0){$+$}}%
    \put(300,68){\line(1,0){70}}%
    \qbezier(300,68)(278,68)(278,48)%
    \qbezier(278,48)(278,28)(300,28)%
    \put(300,28){\line(1,0){70}}%
    \arrL{426}{68}{54}%
    \arrL{345}{56}{40}%
    \arrL{382}{56}{34}%
    \put(298,44){\makebox(0,0){$P$}}%
    \dvh{326}{30}{4}%
    \dvh{356}{30}{4}%
    \arrL{430}{28}{56}%
    \put(444,28){\makebox(0,0){$P$}}%
    \put(335,84){\makebox(0,0){$P{-}Q$}}%
    \put(330,10){\makebox(0,0){(c)}}%
    \put(448,60){\makebox(0,0){$+$}}%
    \put(490,70){\line(1,0){60}}%
    \qbezier(490,70)(478,70)(478,58)%
    \qbezier(478,58)(478,46)(490,46)%
    \qbezier(550,70)(562,70)(562,58)%
    \qbezier(562,58)(562,46)(550,46)%
    \arrL{550}{46}{48}%
    \dvh{508}{24}{5}%
    \dvh{538}{24}{5}%
    \arrL{594}{22}{50}%
    \arrL{642}{22}{46}%
    \put(534,36){\makebox(0,0){$P$}}%
    \put(510,6){\makebox(0,0){(d)}}%
  \end{picture}
\end{equation}
图(21.1a—b)包括一级图(14.10a)，(14.9a)和二级图(14.10b—c)，(14.9b—c)；后者与前者的区别，只在于对内实线的修正；这些实线在图(21.1a—b)中是以粗线表示的，因此它们所对应的不应是理想气体格林函数 $G^{(0)}$，而应是修正到一级项的函数 $G$.最后，(21.1c—d)是(14.10d—e)的二级图.所有的图经过变形，使其结构性质变得更为明显；这就是四外线图的“梯形”级数之前几项，在四外线图内有一对外线按不同方式相互“短路”.

现在我们开始计算图(21.1a).它的解析式为：
\begin{equation}
\begin{gathered}
  \bigl[-\mathrm{i}\Sigma(P)\bigr]_a=\int U(Q)G(P-Q)\frac{\mathrm{d}^4Q}{(2\pi)^4},\\
  Q=(q_0,\bm{q}),\ P=(\omega,\bm{p})
\end{gathered}
\end{equation}
(其中省略了公共因子 $\delta_{\alpha\beta}$).我们首先对 $\mathrm{d}q_0$ 求积分.但是，由于因子 $U(Q)\equiv U(\bm{q})$ 与 $q_0$ 无关，而当 $|q_0|\to\infty$ 时 $G\propto1/q_0$，所以必须预先明确积分方法.为此，需要回顾图(21.1a)的来源，并注意图中的实线对应于同一个算符 $\hat{V}$ 内一对 $\psi$ 算符的收缩.这就是说，$\hat{\Psi}$ 和 $\hat{\Psi}^{+}$ 取在同一时刻，并且当收缩时 $\hat{\Psi}^{+}$ 位于 $\hat{\Psi}$ 的左边.换言之，在坐标表象中产生的 $G$ 函数，取在 $t=t_1-t_2\to-0$.在动量表象中，这意味着对(21.2)式中的被积式添加一个在 $t\to-0$ 时取极限的因子 $\exp(-\mathrm{i}q_0t)$.现在利用公式(7.20)，则得：
\begin{equation}
  \bigl[-\mathrm{i}\Sigma\bigr]_a=\mathrm{i}\int U(\bm{q})N(\bm{p}-\bm{q})\frac{\mathrm{d}^3q}{(2\pi)^3},
\end{equation}
式中 $N(\bm p)$ 为粒子的分布函数.

傅里叶分量 $U(\bm{q})$，只当 $q\geqslant1/r_0$ 时才显著地依赖于 $\bm q$ 的大小，这里 $r_0$ 为场 $U(\bm r)$ 的作用半径；这些 $q$ 值(对于稀薄气体来说)显然比 $p_{\mathrm{F}}$ 大得多，如果限于数值 $|p-p_{\mathrm{F}}|\ll1/r_0$，则对于这些 $\bm q$ 值将有 $N(\bm{p}-\bm{q})\approx0$.因此(21.3)式中的 $U(\bm{q})$ 可以换成 $U(0)$，并从积分号下提出来\footnote{于是不难看出，允许误差的相对数量级 $\sim(p_{\mathrm{F}}r_0)^2$，因此就连对 $p_{\mathrm{F}}r_0$ 下一个数量级的各项也影响不到.}.剩下的积分等于气体密度 $n(\mu)$ 的一半(给定自旋投影值！)，因此 $[\Sigma]_a=-n(\mu)U(0)/2$.

实线自身封闭的图(21.1b)给出了 $[\Sigma]_b=n(\mu)U(0)$.所以，两个图对 $\Sigma$ 的贡献是：
\begin{equation}
  \bigl[\Sigma\bigr]_{a,b}=\frac{1}{2}n(\mu)U(0)=\frac{2\pi}{m}n(\mu)a,
\end{equation}
式中 $a$ 为根据(6.2)式定义的散射长度.

(21.4)式中也含有全部一级效应.在这种近似下，$n(\mu)$ 应理解为理想气体密度 $n^{(0)}(\mu)$，因此，
\begin{equation}
  \Sigma^{(1)}\equiv\bigl[\Sigma\bigr]^{(1)}_{a,b}=\frac{2\pi}{m}n^{(0)}(\mu)a.
\end{equation}
为了向前计算，我们引入作为辅助记号的函数 $F$，它由下列梯形图定义：
\begin{equation}
  \setlength{\unitlength}{0.72pt}%
  \begin{picture}(620,80)(0,0)
    \arrL{98}{56}{46}%
    \arrL{162}{56}{44}%
    \arrL{98}{22}{46}%
    \arrL{162}{22}{44}%
    \put(40,66){\makebox(0,0){$P_3\gamma$}}%
    \put(172,66){\makebox(0,0){$P_1\alpha$}}%
    \put(40,12){\makebox(0,0){$P_4\delta$}}%
    \put(172,12){\makebox(0,0){$P_2\beta$}}%
    \qbezier(130,17)(150.108,18.892)(152,39)%
    \qbezier(152,39)(150.108,59.108)(130,61)%
    \qbezier(130,61)(109.892,59.108)(108,39)%
    \qbezier(108,39)(109.892,18.892)(130,17)%
    \put(130,39){\makebox(0,0){$\mathrm{i}F$}}%
    \put(204,39){\makebox(0,0){$=$}}%
    \arrL{320}{56}{70}%
    \arrL{320}{22}{70}%
    \dvh{285}{28}{4}%
    \put(340,39){\makebox(0,0){$+$}}%
    \arrL{470}{56}{70}%
    \arrL{470}{22}{70}%
    \dvh{425}{28}{4}%
    \dvh{445}{28}{4}%
  \end{picture}
\end{equation}
(照例，$P_1+P_2=P_3+P_4$).其解析形式为：
\begin{equation}
  \mathrm{i}F_{\gamma\delta,\alpha\beta}(P_3,P_4;P_1,P_2)=\mathrm{i}\delta_{\alpha\gamma}\delta_{\beta\delta}\bigl(F^{(1)}+F^{(2)}\bigr),
\end{equation}
其中
\begin{equation}
  \mathrm{i}F^{(1)}=-\mathrm{i}U(P_3-P_1),
\end{equation}
\begin{equation}
  \mathrm{i}F^{(2)}=\int G^{(0)}(P')U(P_1-P')G^{(0)}(P_1+P_2-P')U(P'-P_3)\frac{\mathrm{d}^4P'}{(2\pi)^4}.
\end{equation}
将两个图(21.1c—d)展开，并用 $F^{(2)}$ 把它们表达出来，可得：
\begin{equation}
\begin{aligned}
  \bigl[-\mathrm{i}\Sigma(P)\bigr]_{c,d}={}&-\int G^{(0)}(Q)F^{(2)}(P,Q;Q,P)\frac{\mathrm{d}^4Q}{(2\pi)^4}+\\
  +{}&2\int G^{(0)}(Q)F^{(2)}(P,Q;P,Q)\frac{\mathrm{d}^4Q}{(2\pi)^4}
\end{aligned}
\end{equation}
［其中以 $F^{(1)}$ 代替 $F^{(2)}$ 后，由这两个积分可得出(21.5)式］.两个积分前符号的区别与图(21.1d)中封闭圈的存在有关；第一图中的 $\delta$ 因子给出 $\delta_{\alpha\gamma}\delta_{\gamma\beta}=\delta_{\alpha\beta}$，而第二图中的 $\delta$ 因子给出 $\delta_{\alpha\beta}\delta_{\gamma\gamma}=2\delta_{\alpha\beta}$.

现在来计算 $F^{(2)}$.由于 $U(Q)$ 与 $q_0$ 无关，所以对 $\mathrm{d}p'_0$ 求积分归结为
\begin{equation*}
  \int_{-\infty}^{\infty}G^{(0)}(P')G^{(0)}(P_1+P_2-P')\frac{\mathrm{d}p'_0}{2\pi}.
\end{equation*}
将(9.9)式 $G^{(0)}$ 代入这里(并考虑到 $|p'_0|\to\infty$ 时积分是收敛的)，我们在复数 $p'_0$ 的半平面上用无穷大的半圆圈来封闭积分围道；这时，只当两个函数 $G^{(0)}$ 的极点位于不同的半平面上，积分才不为零，即
\begin{equation}
  \operatorname{sign}(p'-p_{\mathrm{F}})=\operatorname{sign}\bigl(\bigl|\bm{p}_1+\bm{p}_2-\bm{p}'\bigr|-p_{\mathrm{F}}\bigr),
\end{equation}
最后我们得到：
\begin{equation}
\begin{aligned}
  F^{(2)}(P_3,&P_4;P_1,P_2)=\\
  ={}&-\int\frac{U(\bm{p}_1-\bm{p}')U(\bm{p}'-\bm{p}_3)\operatorname{sign}(p'-p_{\mathrm{F}})}
  {\omega_1+\omega_2+2\mu-\dfrac{1}{2m}\bigl[\bm{p}'^2+(\bm{p}_1+\bm{p}_2-\bm{p}')^2\bigr]+\mathrm{i}0\cdot\operatorname{sign}(p'-p_{\mathrm{F}})}
  \frac{\mathrm{d}^3p'}{(2\pi)^3},
\end{aligned}
\end{equation}
(式中 $\omega_1\equiv p_{10},\omega_2\equiv p_{20}$).同时，为了自然地考虑到(21.11)式的要求，在被积式的分子中应作如下代换：
\begin{equation*}
  \operatorname{sign}(p'-p_{\mathrm{F}})\to1-\theta(\bm{p}')-\theta(\bm{p}_1+\bm{p}_2-\bm{p}')
\end{equation*}
其中 $\theta(\bm p)$ 为阶跃函数(1.10).

我们在 §16 中已经看到，梯形图级数确定(真空中)双粒子相互散射幅.因此(21.12)式包含对散射幅各一级项的修正.将(21.8)式的 $F^{(1)}$ 作如下代换后，就能计及这个修正：
\begin{equation*}
  U(\bm{p}_3-\bm{p}_1)\to-\frac{4\pi}{m}\operatorname{Re}f(\bm{p}_3,\bm{p}_1)
\end{equation*}
(其中 $f$ 为精确到二级的真空散射幅)\footnote{切勿将本节中的 $f$ 与准粒子相互作用函数相混淆！}，同时从 $F^{(2)}$ 的表达式(21.12)中减去它在真空中数值的实部，即当 $p_{\mathrm{F}}=0,\mu=0$ 时，数值 $\omega_1=\bm{p}_1^2/2m,\omega_2=\bm{p}_2^2/2m$，它们分别相应于两个真实的碰撞粒子的能量(图的“物理”外线).此后便可以用能量为零时的值，即用散射长度 $a$，去代换 $-\operatorname{Re}f$\footnote{在(21.12)式中不能进行这种代换，因为 $\bm{p}'$ 很大时，将使积分发散. 但施行了上述减法之后(当 $p'\sim p_{\mathrm{F}}$ 时)积分做该代换也已收敛，因而可以进行这种代换. 只减去积分的实部(相应地将 $U$ 换成 $\operatorname{Re}f$)是为了避开与散射幅虚部有关的困难. 因为当动量很小时 $\operatorname{Re}f$ 按动量的偶次幂展开，而 $\operatorname{Im}f$ 按动量的奇次幂展开(见第三卷 §132). 所以考虑 $f$ 的动量依赖关系时，将得出相对数量级为 $(q_{\mathrm{F}}a)^2$ 的修正，即这个修正是可以忽略的. 而将 $U$ 换成 $-4\pi f/m$，则需考虑 $f$ 的虚部，这时将得出相对数量级为 $p_{\mathrm{F}}a$ 的修正.}.所以将有：
\begin{equation}
\begin{aligned}
  F^{(2)}(P_3,&P_4;P_1,P_2)=\\
  ={}&-\biggl(\frac{4\pi a}{m}\biggr)^{\!2}\int\biggl\{\frac{[1-\theta(\bm{p}')-\theta(\bm{p}_1+\bm{p}_2+\bm{p}')]}
  {\omega_1+\omega_2+2\mu-\dfrac{1}{2m}\bigl[\bm{p}'^2+(\bm{p}_1+\bm{p}_2-\bm{p}')^2\bigr]+\mathrm{i}0\cdot\operatorname{sign}(p'-p_{\mathrm{F}})}-\\
  -\mathrm{P}&\,\frac{2m}{p_1^2+p_2^2-\bm{p}'^2-(\bm{p}_1+\bm{p}_2-\bm{p}')^2}\biggr\}\frac{\mathrm{d}^3p'}{(2\pi)^3}.
\end{aligned}
\end{equation}
第二项中的记号 P，表示取主值积分；这是利用规则(8.11)分离出积分实部的结果.

因为(21.13)式对 $P_1$ 和 $P_2$ 是对称的，(21.10)式中的两个积分相等，因此
\begin{equation*}
  \bigl[-\mathrm{i}\Sigma(P)\bigr]_{c,d}=\int G^{(0)}(Q)F^{(2)}(P,Q;P,Q)\frac{\mathrm{d}^4Q}{(2\pi)^4}.
\end{equation*}
将(21.13)式的第一项代入上式时，如果
\begin{equation}
  \operatorname{sign}(p'-p_{\mathrm{F}})=-\operatorname{sign}(q-p_{\mathrm{F}}),
\end{equation}
则对 $\mathrm{d}q_0$ 的积分不等于零，因此被积式的两个极点重新处于 $q_0$ 的不同的半平面上.将(21.13)式的第二项代入时，只有因子 $G_0(Q)$ 与 $q_0$ 有关，利用公式(7.23)可求出对 $\mathrm{d}q_0$ 的积分，并给出 $N^{(0)}(\bm{q})$——理想气体粒子的分布函数，即阶跃函数 $\theta(\bm q)$.将(21.1a—d)全部图的贡献汇聚起来，结果我们得到：
\begin{equation}
  \Sigma(\omega,\bm{p})=\frac{2\pi}{m}n(\mu)a+\Sigma^{(2)}(\omega,\bm{p}),
\end{equation}
式中
\begin{equation}
\begin{aligned}
  \Sigma^{(2)}(\omega,&\bm{p})=\biggl(\frac{4\pi a}{m}\biggr)^{\!2}\int\biggl\{\frac{[1-\theta(\bm{p}')-\theta(\bm{p}+\bm{q}-\bm{p}')]\bigl[\theta(\bm{q})-\theta(\bm{p}')\bigr]}
  {\omega+\mu+\dfrac{1}{2m}\bigl[q^2-\bm{p}'^2-(\bm{p}+\bm{q}-\bm{p}')^2\bigr]+\mathrm{i}0\cdot\operatorname{sign}(p'-p_{\mathrm{F}})}-\\
  -\mathrm{P}&\,\frac{2m\theta(\bm{q})}{p^2+q^2-\bm{p}'^2+(\bm{p}+\bm{q}-\bm{p}')^2}\biggr\}\frac{\mathrm{d}^3q\,\mathrm{d}^3p'}{(2\pi)^6}.
\end{aligned}
\end{equation}
［积分号下第一项分子中的因子 $\theta(\bm{q})-\theta(\bm{p}')$，当满足条件(21.14)时则换成 $-\operatorname{sign}(q-p_{\mathrm{F}})$］.

我们首先指出，$\Sigma$ 具有虚部.这个虚部可以借助规则(8.11)从(21.16)式中分解出来，并用下式表出：
\begin{equation}
\begin{aligned}
  \operatorname{Im}\Sigma(\omega,&\bm{p})=-\biggl(\frac{4\pi a}{m}\biggr)^{\!2}\pi\int\bigl\{\theta(\bm{q})\bigl[1-\theta(\bm{p}')\bigr]\bigl[1-\theta(\bm{p}+\bm{q}-\bm{p}')\bigr]-\\
  -{}&\bigl[1-\theta(\bm{q})\bigr]\theta(\bm{p}')\theta(\bm{p}+\bm{q}-\bm{p}')\bigr\}\times\\
  \times{}&\delta\biggl[\omega+\mu+\frac{1}{2m}\bigl(q^2-\bm{p}'^2-(\bm{p}+\bm{q}-\bm{p}')^2\bigr)\biggr]\frac{\mathrm{d}^3q\,\mathrm{d}^3p'}{(2\pi)^6}
\end{aligned}
\end{equation}
［考虑到 $\theta^2(\bm{p})\equiv\theta(\bm{p})$，已对花括号中的表达式作了变换］.

准粒子能谱，可根据(14.13)式计算出来，即
\begin{equation}
  \varepsilon(\bm{p})=\frac{p^2}{2m}+\frac{2\pi}{m}n(\mu)a+\Sigma^{(2)}\biggl(\frac{p^2}{2m}-\mu,\bm{p}\biggr)
\end{equation}
(在 $\Sigma^{(2)}$ 中，按需要的精确度可取 $\varepsilon\approx p^2/2m$).$\Sigma$ 的复数性表示激发有衰减($\operatorname{Im}\varepsilon\neq0$).

出现这种衰减，表明准粒子是不稳定的，这种不稳定性与准粒子可能有实际裂变过程有关.准粒子可以放出自己的部分能量，并靠它产生准粒子对(粒子和空穴).作为例子，我们来考查(21.17)式积分号下花括号中的第一项.根据阶跃函数的性质，如果
\begin{equation*}
  p'>p_{\mathrm{F}},\quad\bigl|\bm{q}+\bm{p}-\bm{p}'\bigr|>p_{\mathrm{F}},\quad q<p_{\mathrm{F}},
\end{equation*}
则这一项不等于零.这几个不等式对应于这样一种过程：在此过程中，初动量为 $\bm{p}$($p>p_{\mathrm{F}}$)的准粒子转变为动量为 $\bm{p}'$ 的状态($p>p'>p_{\mathrm{F}}$)，并且有 $\bm{p}-\bm{p}'$ 的动量传递给费米球内(动量 $q<p_{\mathrm{F}}$)的粒子，后者被激发到费米球外动量为 $\bm{q}+\bm{p}-\bm{p}'$ 的状态；这种转变相当于出现了动量为 $-\bm{q}$(空穴)和 $\bm{q}+\bm{p}-\bm{p}'$ 的两个新的元激发.(21.17)式中的 $\delta$ 函数，表明在这个过程中遵从能量守恒定律，其中 $\omega+\mu$ 起着准粒子初始能量 $\varepsilon(\bm{p})$ 的作用：
\begin{equation*}
  \varepsilon(\bm{p})=\varepsilon(\bm{p}')+\bigl[\varepsilon(\bm{q}+\bm{p}-\bm{p}')-\varepsilon(\bm{q})\bigr]
\end{equation*}
(此处，在一级近似下，只要取 $\varepsilon(\bm{p})=\bm{p}^2/2m$ 就够了).按照上面指出的意义，以该等式确定的能量 $\varepsilon(\bm p)$，确实相应于费米球外($\varepsilon>\mu$)的准粒子.

与此类似，(21.17)式花括号内的第二项，是准粒子对产生于空穴的过程.此项给出了能量为 $\varepsilon<\mu$ 的元激发的衰减.用图技术语言来说，准粒子可能产生对，表现为可能通过截断三条实线(其中，两条同方向，第三条反方向)而将 $G$ 函数图分成两部分.在图(21.1c—d)中，这种分割是在两条虚线之间进行的.

弱非理想气体情况是特殊的(与任意费米液体的一般情况相比)，原因在于，这种情况中的准粒子能谱，不仅在费米面附近而且在整个动量值域内都有意义：因为“气态参数” $ap_{\mathrm{F}}$ 本是一个很小的量，所以准粒子衰减($\operatorname{Im}\varepsilon$)相当小.但是，这里我们只对两种极端情况得出最终的计算结果.

在费米面附近($|p-p_{\mathrm{F}}|\ll p_{\mathrm{F}}$)，得到：
\begin{equation*}
  \operatorname{Re}\varepsilon=\mu+(p-p_{\mathrm{F}})p_{\mathrm{F}}/m^{*},
\end{equation*}
其中 $\mu$ 和 $m^{*}$ 分别由(6.14)和(6.17)式决定.得出的准粒子衰减为：
\begin{equation}
  \operatorname{Im}\varepsilon=-\frac{1}{\pi m}(p_{\mathrm{F}}a)^2(p-p_{\mathrm{F}})^2\operatorname{sign}(p-p_{\mathrm{F}}).
\end{equation}
此表达式与 $(p-p_{\mathrm{F}})^2$ 成比例具有明显的来源：所出现的一个因子 $p-p_{\mathrm{F}}$，乃是动量空间一个区域(薄球壳)的厚度，准粒子对所由产生的准粒子，其动量即属于此区域.还有一个同样的因子，则是在其中产生准粒子对的壳层的厚度.我们顺便指出，这些见解对任何费米液体都是适用的，因此在费米面附近总有 $\operatorname{Im}\varepsilon\propto(p-p_{\mathrm{F}})^2$\footnote{当温度不等于零时，将这个量对热平衡分布求平均，可得出准粒子衰减与 $T^2$ 成比例的结果，关于这一点已在 §1 中讲过.}.

在大动量的情况下 $p\gg p_{\mathrm{F}}$(但仍保持 $pa\ll1$)，我们有：
\begin{equation}
  \varepsilon=\biggl(\frac{p^2}{2m}+\frac{2p_{\mathrm{F}}^2}{3\pi m}p_{\mathrm{F}}a\biggr)-\mathrm{i}\,\frac{p_{\mathrm{F}}p}{3\pi m}\bigl(p_{\mathrm{F}}a\bigr)^2.
\end{equation}
在两种情况下，比值 $\operatorname{Im}\varepsilon/\operatorname{Re}\varepsilon$ 都小.当 $p\sim p_{\mathrm{F}}$ 时，这个比值达到最大值，即使在这里它 $\sim(p_{\mathrm{F}}a)^2\ll1$.

最后，我们引入弱非理想气体格林函数的重整化常数值.该常数可用如下公式算出：
\begin{equation*}
  \frac{1}{Z}=1-\frac{\partial\Sigma(\omega,\bm{p})}{\partial\omega}\bigg|_{\omega=0,\,p=p_{\mathrm{F}}},
\end{equation*}
因而它等于
\begin{equation}
  Z=1-\frac{8\ln 2}{\pi^2}\bigl(p_{\mathrm{F}}a\bigr)^2.
\end{equation}
